\chapter{Maximum Principle and Li-Yau-Hamilton Inequalities}
\label{chap:cz-maximum-principles}

\bigskip
The maximum principle is a fundamental tool in the study of
parabolic equations in general. In this chapter, we present
various maximum principles for tensors developed by Hamilton in
the Ricci flow. As an immediate consequence, the Ricci flow
preserves the nonnegativity of the curvature operator. We also
present two crucial estimates in the Ricci flow: the Hamilton-Ivey
curvature pinching estimate \cite{Ha95F, Iv} when dimension $n=3$,
and the Li-Yau-Hamilton estimate \cite{Ha88, Ha93} from which one
obtains the Harnack inequality for the evolved scalar curvature
via a Li-Yau \cite{LY} path integral. Most of the presentation in
Sections 2.1-2.5 follows closely Hamilton \cite{Ha82, Ha86, Ha88,
Ha93, Ha95F}, and some parts of Section 2.3 also follows Chow-Lu
\cite{Cl04}. Finally, in Section 2.6 we describe Perelman's Li-Yau
type estimate for solutions to the conjugate heat equation and
show how Li-Yau type path integral leads to a space-time distance
function (i.e., what Perelman \cite{P1} called the reduced
distance).


\section{Preserving Positive Curvature}

Let $M$ be an $n$-dimensional complete manifold. Consider a family
of smooth metrics $g_{ij}(t)$ evolving by the Ricci flow with
uniformly bounded curvature for $t\in[0,T]$ with $T<+\infty$.
Denote by $d_t(x,y)$ the distance between two points $x,y\in M$
with respect to the metric $g_{ij}(t)$. First we need the
following useful fact (cf. \cite{ScY})

%%\vskip 0.2cm\noindent {\bf Lemma 2.1.1}\ \emph{
% CZ-BASELINE source-line:3030 source-kind:lemma source-id:2.1.1
\begin{lemma}[proper exhaustion with uniformly bounded derivatives]
  \label{lem:cz-proper-exhaustion-with-uniformly-bounded-derivatives}
  \dcref{cao-zhu:2.1.1}
  \group{Cao--Zhu Preserving Positive Curvature}
  \source{cao-zhu}
There exists a smooth function $f$ on $M$ such that $f\ge 1$
everywhere, $f(x)\rightarrow +\infty$ as $d_0(x,x_0)\rightarrow
+\infty$ $($for some fixed $x_0\in M),$
$$
|\nabla f|_{g_{ij}(t)}\le C \quad \text{and}\quad |\nabla^2
f|_{g_{ij}(t)}\le C
$$
on $M \times [0,T]$ for some positive constant $C$.
\end{lemma}

%%\vskip 0.1cm\noindent {\bf Proof.}\
\begin{proof}
  \uses{thm:cz-global-curvature-derivative-estimates}
Let $\varphi(v)$ be a smooth function on $\mathbb{R}^n$ which is
nonnegative, rotationally symmetric and has compact support in a
small ball centered at the origin with
$\int_{\mathbb{R}^n}\varphi(v)dv=1$.

For each $x\in M$, set
$$
f(x)=\int_{\mathbb{R}^n}\varphi(v)(d_0(x_0,exp_x(v))+1)dv,
$$
where the integral is taken over the tangent space $T_xM$ at $x$
which we have identified with $\mathbb{R}^n$. If the size of the
support of $\varphi(v)$ is small compared to the maximum
curvature, then it is well known that this defines a smooth
function $f$ on $M$ with $f(x)\rightarrow +\infty$ as
$d_0(x,x_0)\rightarrow +\infty$, while the bounds on the first and
second covariant derivatives of $f$ with respect to the metric
$g_{ij}(\cdot,0)$ follow from the Hessian comparison theorem. Thus
it remains to show these bounds hold with respect to the evolving
metric ${g_{ij}(t)}$.

We compute, using the frame $\{F_a^i\nabla_if\}$ introduced in
Section 1.3,
$$
\frac{\partial}{\partial t}\nabla_a f =\frac{\partial}{\partial
t}(F_a^i\nabla_if) =R_{ab}\nabla_b f.
$$
Hence
$$
|\nabla f|\le C_1\cdot e^{C_2 t},
$$
where $C_1, C_2$ are some positive constants depending only on the
dimension. Also
\begin{align*}
\frac{\partial}{\partial t}(\nabla_a\nabla_bf) &
=\frac{\partial}{\partial t} \(F_a^iF_b^j\(\frac{\partial^2
f}{\partial x^i\partial x^j}
-\Gamma_{ij}^k\frac{\partial f}{\partial x^k}\)\)\\
& =R_{ac}\nabla_b\nabla_cf+R_{bc}\nabla_a\nabla_cf
+(\nabla_cR_{ab}-\nabla_aR_{bc}-\nabla_bR_{ac})\nabla_cf.
\end{align*}
Then by Shi's derivative estimate (Theorem 1.4.1), we have
$$
\frac{\partial}{\partial t}|\nabla^2 f|\le C_3|\nabla^2 f|
+\frac{C_3}{\sqrt{t}},
$$
which implies
$$
|\nabla^2f|_{g_{ij}(t)}\le e^{C_3 t}\(|\nabla^2f|_{g_{ij}(0)}
+\int_0^t\frac{C_3}{\sqrt{\tau}}e^{-C_3\tau}d\tau\)
$$
for some positive constants $C_3$ depending only on the dimension
and the curvature bound.
%%$$\eqno \#$$ \vskip 0.3cm
\end{proof}

We now use the weak maximum principle to derive the following result
(cf. \cite{Ha82} and \cite{Sh89}).

%%\vskip 0.2cm\noindent {\bf Proposition 2.1.2}\ \emph{
% CZ-BASELINE source-line:3102 source-kind:proposition source-id:2.1.2
\begin{proposition}[preservation of nonnegative scalar curvature]
  \label{prop:cz-preservation-of-nonnegative-scalar-curvature}
  \dcref{cao-zhu:2.1.2}
  \group{Cao--Zhu Preserving Positive Curvature}
  \source{cao-zhu}
If the scalar curvature $R$ of the solution $g_{ij}(t), 0\le t\le
T$, to the Ricci flow is nonnegative at $t=0$, then it remains so
on $0\le t\le T$.
\end{proposition}

%%\vskip 0.1cm\noindent {\bf Proof.}\
\begin{proof}
  \uses{lem:cz-proper-exhaustion-with-uniformly-bounded-derivatives}
Let $f$ be the function constructed in Lemma 2.1.1 and recall
$$
\frac{\partial R}{\partial t}=\Delta R+2|\Ric|^2.
$$
For any small constant $\varepsilon>0$ and large constant $A>0$,
we have
\begin{align*}
\frac{\partial}{\partial t}(R+\varepsilon e^{At}f)
&  = \frac{\partial R}{\partial t}+\varepsilon Ae^{At}f\\
&  = \Delta(R+\varepsilon e^{At}f)+2|Ric|^2+\varepsilon
e^{At}(Af-\Delta f)\\
& > \Delta(R+\varepsilon e^{At}f)
\end{align*}
by choosing $A$ large enough.

We claim that
$$
R+\varepsilon e^{At}f>0\qquad \text{on}\quad M\times[0,T].
$$
Suppose not, then there exist a first time $t_0>0$ and a point
$x_0\in M$ such that
\begin{align*}
(R+\varepsilon e^{At}f)(x_0,t_0)&=0,\\
\nabla(R+\varepsilon e^{At}f)(x_0,t_0)&=0,\\
\Delta(R+\varepsilon e^{At}f)(x_0,t_0)&\ge 0,\\
\text{and}\quad \frac{\partial}{\partial t}(R+\varepsilon
e^{At}f)(x_0,t_0)&\le 0.
\end{align*}
Then
$$
0\ge \frac{\partial}{\partial t}(R+\varepsilon
e^{At}f)(x_0,t_0)>\Delta(R+\varepsilon e^{At}f)(x_0,t_0)\ge 0,
$$
which is a contradiction. So we have proved that
$$
R+\varepsilon e^{At}f>0\qquad \text{on}\quad M\times[0,T].
$$
Letting $\varepsilon\rightarrow 0$, we get
$$
R\ge 0\qquad \text{on}\quad M\times[0,T].
$$
This finishes the proof of the proposition.
%%$$\eqno \#$$\vskip 0.3cm
\end{proof}

Next we derive a maximum principle of Hamilton for tensors. Let
$M$ be a complete manifold with a metric $g=\{g_{ij}\}$, $V$ a
vector bundle over $M$ with a metric $h=\{h_{\alpha\beta}\}$ and a
connection $\nabla=\{\Gamma^\alpha_{i\beta}\}$ compatible with
$h$, and suppose $h$ is fixed but $g$ and $\nabla$ may vary
smoothly with time $t$. Let $\Gamma(V)$ be the vector space of
$C^\infty$ sections of $V$. The Laplacian $\Delta$ acting on a
section $\sigma\in\Gamma(V)$ is defined by
$$
\Delta\sigma=g^{ij}\nabla_i\nabla_j\sigma.
$$
Let $M_{\alpha\beta}$ be a symmetric bilinear form on $V$. We say
$M_{\alpha\beta}\ge 0$ if $M_{\alpha\beta}v^\alpha v^\beta\ge 0$
for all vectors $v=\{v^\alpha\}$. Assume
$N_{\alpha\beta}={\mathcal{P}}(M_{\alpha\beta},h_{\alpha\beta})$
is a polynomial in $M_{\alpha\beta}$ formed by contracting
products of $M_{\alpha\beta}$ with itself using the metric
$h=\{h_{\alpha\beta}\}$. Assume that the tensor $M_{\alpha\beta}$
is uniformly bounded in space-time and let $g_{ij}$ evolve by the
Ricci flow with bounded curvature.

%%\vskip 0.2cm\noindent{\bf Lemma 2.1.3} (Hamilton \cite{Ha82})\emph{
% CZ-BASELINE source-line:3177 source-kind:lemma source-id:2.1.3
\begin{lemma}[tensor maximum principle on complete manifolds]
  \label{lem:cz-tensor-maximum-principle-on-complete-manifolds}
  \dcref{cao-zhu:2.1.3}
  \group{Cao--Zhu Preserving Positive Curvature}
  \source{cao-zhu}
Suppose that on $0\le t\le T$,
$$
\frac{\partial}{\partial t}M_{\alpha\beta} =\Delta
M_{\alpha\beta}+u^i\nabla_i M_{\alpha\beta}+N_{\alpha\beta}
$$
where $u^i(t)$ is a time-dependent vector field on $M$ with
uniform bound and
$N_{\alpha\beta}={\mathcal{P}}(M_{\alpha\beta},h_{\alpha\beta})$
satisfies
$$
N_{\alpha\beta}v^\alpha v^\beta\ge 0 \quad \text{whenever} \quad
M_{\alpha\beta}v^\beta=0.
$$
If $M_{\alpha\beta}\ge 0$ at $t=0$, then it remains so on $0\le
t\le T$.
\end{lemma}

%%\vskip 0.1cm\noindent{\bf Proof.}\
\begin{proof}
  \uses{lem:cz-proper-exhaustion-with-uniformly-bounded-derivatives}
Set
$$
\tilde{M}_{\alpha\beta} =M_{\alpha\beta}+\varepsilon e^{At}f
h_{\alpha\beta},
$$
where $A>0$ is a suitably large constant (to be chosen later) and
$f$ is the function constructed in Lemma 2.1.1.

We claim that $\tilde{M}_{\alpha\beta}>0$ on $M\times[0,T]$ for
every $\varepsilon>0$. If not, then for some $\varepsilon>0$,
there will be a first time $t_0>0$ where $\tilde{M}_{\alpha\beta}$
acquires a null vector $v^\alpha$ of unit length at some point
$x_0\in M$. At $(x_0,t_0)$,
\begin{align*}
N_{\alpha\beta}v^\alpha v^\beta & \geq N_{\alpha\beta}v^\alpha
v^\beta
-\tilde{N}_{\alpha\beta}v^\alpha v^\beta\\
& \ge -C\varepsilon e^{At_0}f(x_0),
\end{align*}
where
$\tilde{N}_{\alpha\beta}=\mathcal{P}(\tilde{M}_{\alpha\beta},
h_{\alpha\beta})$, and $C$ is a positive constant (depending on
the bound of $M_{\alpha\beta}$, but independent of $A$).

Let us extend $v^\alpha$ to a local vector field in a neighborhood
of $x_0$ by parallel translating $v^{\alpha}$ along geodesics
(with respect to the metric $g_{ij}(t_0)$) emanating radially out
of $x_0$, with $v^{\alpha}$ independent of $t$. Then, at
$(x_0,t_0)$, we have
\begin{align*}
\frac{\partial}{\partial t}(\tilde{M}_{\alpha\beta}v^\alpha
v^\beta)&  \le 0 ,\\
\nabla (\tilde{M}_{\alpha\beta}v^\alpha v^\beta)&  = 0, \\
 \text{and} \quad \Delta(\tilde{M}_{\alpha\beta}v^\alpha
v^\beta)&  \ge 0.
\end{align*}
But
\begin{align*}
0\ge\frac{\partial}{\partial t}(\tilde{M}_{\alpha\beta}v^\alpha
v^\beta) &  = \frac{\partial}{\partial t}(M_{\alpha\beta}v^\alpha
v^\beta+\varepsilon e^{At}f),\\
&  = \Delta(\tilde{M}_{\alpha\beta}v^\alpha
v^\beta)-\Delta(\varepsilon
e^{At}f)+u^i\nabla_i(\tilde{M}_{\alpha\beta}v^\alpha
v^\beta)\\
&\quad -u^i\nabla_i(\varepsilon e^{At}f)+N_{\alpha\beta}v^\alpha
v^\beta+\varepsilon A e^{At_0}f(x_0)\\
&  \ge -C\varepsilon e^{At_0}f(x_0)+\varepsilon A e^{At_0}f(x_0)>0
\end{align*}
when $A$ is chosen sufficiently large. This is a contradiction.
%$$\eqno \#$$ \vskip 0.3cm
\end{proof}

By applying Lemma 2.1.3 to the evolution equation
$$
\frac{\partial}{\partial t}M_{\alpha\beta} =\Delta
M_{\alpha\beta}+M_{\alpha\beta}^2+M_{\alpha\beta}^\#
$$
of the curvature operator $M_{\alpha\beta}$, we immediately obtain
the following important result.

%%\vskip 0.2cm\noindent{\bf Proposition 2.1.4} (Hamilton
%%\cite{Ha86}) \ \ \emph{
% CZ-BASELINE source-line:3260 source-kind:proposition source-id:2.1.4
\begin{proposition}[preservation of nonnegative curvature operator]
  \label{prop:cz-preservation-of-nonnegative-curvature-operator}
  \dcref{cao-zhu:2.1.4}
  \group{Cao--Zhu Preserving Positive Curvature}
  \source{cao-zhu}
Nonnegativity of the curvature operator $M_{\alpha\beta}$ is
preserved by the Ricci flow.
\end{proposition}   %%$$\eqno \#$$ \vskip 0.3cm

In the K\"ahler case, the nonnegativity of the holomorpic
bisectional curvature is preserved under the K\"ahler-Ricci flow.
This result is proved by Bando \cite{Bando} for complex dimension
$n=3$ and by Mok \cite{Mo88} for general dimension $n$ when the
manifold is compact, and by Shi \cite{Sh90} when the manifold is
noncompact.

%%\vskip 0.2cm\noindent{\bf Proposition 2.1.5} \   \emph{
% CZ-BASELINE source-line:3273 source-kind:proposition source-id:2.1.5
\begin{proposition}[preservation of holomorphic bisectional curvature]
  \label{prop:cz-preservation-of-holomorphic-bisectional-curvature}
  \dcref{cao-zhu:2.1.5}
  \group{Cao--Zhu Preserving Positive Curvature}
  \source{cao-zhu}
Under the K\"ahler-Ricci flow if the initial metric has positive
$($nonnegative$)$ holomorphic bisectional curvature then the
evolved metric also has positive $($nonnegative$)$ holomorphic
bisectional curvature.
\end{proposition}   %%$$\eqno\#$$\vskip 1cm


\section{Strong Maximum Principle}

Let $\Omega$ be a bounded, connected open set of a complete
$n$-dimensional manifold $M$, and let $g_{ij}(x,t)$ be a smooth
solution to the Ricci flow on $\Omega \times [0,T]$. Consider a
vector bundle $V$ over $\Omega$ with a fixed metric
$h_{\alpha\beta}$ (independent of time), and a connection $\nabla
= \{\Gamma^\alpha_{i\beta}\}$ which is compatible with
$h_{\alpha\beta}$ and may vary with time $t$. Let $\Gamma(V)$ be
the vector space of $C^\infty$ sections of $V$ over $\Omega$. The
Laplacian $\Delta$ acting on a section $\sigma\in\Gamma(V)$ is
defined by
$$
\Delta\sigma=g^{ij}(x,t)\nabla_i \nabla_j\sigma.
$$

Consider a family of smooth symmetric bilinear forms
$M_{\alpha\beta}$ evolving by \be \frac{\partial}{\partial
t}M_{\alpha\beta}=\Delta M_{\alpha\beta}+N_{\alpha\beta},\ \
\text{ on }\
\Omega\times[0,T], %\eqno(2.2.1)
\ee where $N_{\alpha\beta}=P(M_{\alpha\beta},h_{\alpha\beta})$ is
a polynomial in $M_{\alpha\beta}$ formed by contracting products
of $M_{\alpha\beta}$ with itself using the metric
$h_{\alpha\beta}$ and satisfies
$$
N_{\alpha\beta}\geq 0, \; \mbox{ whenever }\; M_{\alpha\beta}\geq
0.
$$
The following result, due to Hamilton \cite{Ha86}, shows that the
solution of (2.2.1) satisfies a strong maximum principle.

%%\vskip 0.2cm \noindent {\bf Theorem 2.2.1}
% CZ-BASELINE source-line:3314 source-kind:theorem source-id:2.2.1
\begin{theorem}[strong maximum principle for symmetric bilinear forms]
  \label{thm:cz-strong-maximum-principle-for-symmetric-bilinear-forms}
  \dcref{cao-zhu:2.2.1}
  \group{Cao--Zhu Strong Maximum Principle}
  \source{cao-zhu}
 Let
$M_{\alpha\beta}$ be a smooth solution of the equation $(2.2.1)$.
Suppose $M_{\alpha\beta}\geq0$ on $\Omega \times [0,T]$. Then
there exists a positive constant $0<\delta \leq T$ such that on
$\Omega \times (0,\delta)$, the rank of $M_{\alpha\beta}$ is
constant, and the null space of $M_{\alpha\beta}$ is invariant
under parallel translation and invariant in time and also lies in
the null space of $N_{\alpha\beta}$.
\end{theorem}

%%\vskip 0.1cm \noindent {\bf Proof.}
\begin{proof}
Set
$$
l=\max_{x\in \Omega}\{\text{rank of } M_{\alpha\beta}(x,0)\}.
$$
Then we can find a nonnegative smooth function $\rho(x)$, which is
positive somewhere and has compact support in $\Omega$, so that at
every point $x \in \Omega$,
$$
\sum^{n-l+1}_{i=1}M_{\alpha\beta}(x,0)v^{\alpha}_{i}v^{\beta}_{i}
\geq \rho(x)
$$
for any $(n-l+1)$ orthogonal unit vectors $\{ v_1, \ldots,
v_{n-l+1} \}$ at $x$.

Let us evolve $\rho(x)$ by the heat equation
$$
\frac{\partial}{\partial t}\rho=\Delta\rho
$$
with the Dirichlet condition $\rho|_{\partial \Omega}=0$ to get a
smooth function $\rho(x,t)$ defined on $\Omega \times [0,T]$. By
the standard strong maximum principle, we know that $\rho(x,t)$ is
positive everywhere in $\Omega$ for all $t\in (0,T]$.

For every $\varepsilon>0$, we claim that at every point $(x,t) \in
\Omega \times [0,T]$, there holds
$$
\sum^{n-l+1}_{i=1}M_{\alpha\beta}(x,t)v^{\alpha}_{i}v^{\beta}_{i}
 + \varepsilon e^t > \rho(x,t)
$$
for any $(n-l+1)$ orthogonal unit vectors $\{ v_1, \ldots,
v_{n-l+1} \}$ at $x$.

We argue by contradiction. Suppose not, then for some
$\varepsilon>0$, there will be a first time $t_0 > 0$ and some
$(n-l+1)$ orthogonal unit vectors $\{ v_1, \ldots, v_{n-l+1} \}$
at some point $x_0 \in \Omega$ so that
$$
\sum^{n-l+1}_{i=1}M_{\alpha\beta}(x_0,t_0)v^{\alpha}_{i}v^{\beta}_{i}
 + \varepsilon e^{t_0} = \rho(x_0,t_0)
$$

Let us extend each $v_i$ ($i=1, \ldots, n-l+1$) to a local vector
field, independent of $t$, in a neighborhood of $x_0$ by parallel
translation along geodesics (with respect to the metric
$g_{ij}(t_0)$) emanating radially out of $x_0$. Clearly $\{v_1,
\ldots, v_{n-l+1}\}$ remain orthogonal unit vectors in the
neighborhood.  Then, at $(x_0,t_0)$, we have
\begin{align*}
\frac{\partial}{\partial t}
\(\sum^{n-l+1}_{i=1}M_{\alpha\beta}v^{\alpha}_{i}v^{\beta}_{i} +
\varepsilon e^t -\rho\)
& \le 0 ,\\
\text{and} \quad
\Delta\(\sum^{n-l+1}_{i=1}M_{\alpha\beta}v^{\alpha}_{i}v^{\beta}_{i}
 + \varepsilon e^t -\rho\)&  \ge 0.
\end{align*}
But, since $N_{\alpha\beta}\ge 0$ by our assumption, we have
\begin{align*}
0&  \ge \frac{\partial}{\partial t}
\(\sum^{n-l+1}_{i=1}M_{\alpha\beta}v^{\alpha}_{i}v^{\beta}_{i}
+ \varepsilon e^t -\rho\)\\
&  = \sum^{n-l+1}_{i=1}(\Delta
M_{\alpha\beta}+N_{\alpha\beta})v^{\alpha}_{i}v^{\beta}_{i}
 + \varepsilon e^t -\Delta \rho\\
& \geq\sum^{n-l+1}_{i=1}\Delta(M_{\alpha\beta}
v^{\alpha}_{i}v^{\beta}_{i})+ \varepsilon e^t -\Delta \rho\\
&  =\sum^{n-l+1}_{i=1}\Delta(M_{\alpha\beta}v^{\alpha}_{i}
v^{\beta}_{i}+\varepsilon e^t -\rho)+ \varepsilon e^t \\
& \ge \varepsilon e^t >0.
\end{align*}
This is a contradiction. Thus by letting $\varepsilon\rightarrow
0$, we prove that
$$
\sum^{n-l+1}_{i=1}M_{\alpha\beta}(x,t)v^{\alpha}_{i}v^{\beta}_{i}
 \geq \rho(x,t)
$$
for any $(n-l+1)$ orthogonal unit vectors $\{ v_1, \ldots,
v_{n-l+1} \}$ at $x \in \Omega$ and $t \in [0,T]$. Hence
$M_{\alpha\beta}$ has at least rank $l$ everywhere in the open set
$\Omega$ for all $t \in (0,T]$. Therefore we can find a positive
constant $\delta (\leq T)$ such that the rank $M_{\alpha\beta}$ is
constant over $\Omega \times (0,\delta)$.

Next we proceed to analyze the null space of $M_{\alpha\beta}$.
Let $v$ be any smooth section of $V$ in the null of
$M_{\alpha\beta}$ on $0<t<\delta$. Then
\begin{align*}
0&  = \frac{\partial}{\partial t}(M_{\alpha\beta}v^\alpha
v^\beta)\\
&  = \(\frac{\partial}{\partial t}M_{\alpha\beta}\)v^\alpha
v^\beta+2M_{\alpha\beta}v^\alpha\frac{\partial v^\beta}{\partial t}\\
&  = \(\frac{\partial}{\partial t}M_{\alpha\beta}\)v^\alpha
v^\beta,
\end{align*}
and
\begin{align*}
0&  = \Delta(M_{\alpha\beta}v^\alpha v^\beta)\\
&  = (\Delta M_{\alpha\beta})v^\alpha v^\beta
+4g^{kl}\nabla_kM_{\alpha\beta}\cdot v^\alpha\nabla_lv^\beta\\
&\quad +2M_{\alpha\beta}g^{kl}\nabla_kv^\alpha\cdot\nabla_lv^\beta
+2M_{\alpha\beta}v^\alpha\Delta v^\beta\\
&  = (\Delta M_{\alpha\beta})v^\alpha
v^\beta+4g^{kl}\nabla_kM_{\alpha\beta}\cdot
v^\alpha\nabla_lv^\beta
+2M_{\alpha\beta}g^{kl}\nabla_kv^\alpha\cdot\nabla_lv^\beta.
\end{align*}
By noting that
$$
0=\nabla_k(M_{\alpha\beta}v^\beta)
=(\nabla_kM_{\alpha\beta})v^\alpha+M_{\alpha\beta}\nabla_kv^\alpha
$$
and using the evolution equation (2.2.1), we get
$$
N_{\alpha\beta}v^\alpha v^\beta
+2M_{\alpha\beta}g^{kl}\nabla_kv^\alpha\cdot\nabla_lv^\beta=0.
$$
Since $M_{\alpha\beta}\ge 0$ and $N_{\alpha\beta}\ge 0$, we must
have
$$
v\in \text{null}\; (N_{\alpha\beta})\quad \text{and} \quad
\nabla_i v\in \text{null}\; (M_{\alpha\beta}),\quad \mbox{for all
}\; i.
$$
The first inclusion shows that $\text{null}\;
(M_{\alpha\beta})\subset \text{null}\; (N_{\alpha\beta}),$ and the
second inclusion shows that $\text{null}\; (M_{\alpha\beta})$ is
invariant under parallel translation.

To see $\text{null}\; (M_{\alpha\beta})$ is also invariant in
time, we first note that
$$
\Delta v=\nabla^i(\nabla_i v)\in \text{null}\; (M_{\alpha\beta})
$$
and then
$$
g^{kl}\nabla_kM_{\alpha\beta}\cdot\nabla_lv^\alpha
=g^{kl}\nabla_k(M_{\alpha\beta}\nabla_lv^\alpha)
-M_{\alpha\beta}\Delta v^\alpha=0.
$$
Thus we have
\begin{align*}
0& = \Delta(M_{\alpha\beta}v^\alpha)\\
& = (\Delta M_{\alpha\beta})v^\alpha
+2g^{kl}\nabla_kM_{\alpha\beta}\cdot \nabla_lv^\alpha
+M_{\alpha\beta}\Delta v^\alpha\\
&  = (\Delta M_{\alpha\beta})v^\alpha,
\end{align*}
and hence
\begin{align*}
0&  = \frac{\partial}{\partial t}(M_{\alpha\beta}v^\alpha)\\
&  = (\Delta M_{\alpha\beta}+N_{\alpha\beta})v^\alpha
+M_{\alpha\beta}\frac{\partial v^\alpha}{\partial t}\\
&  = M_{\alpha\beta}\frac{\partial v^\alpha}{\partial t}.
\end{align*}
This shows that
$$
\frac{\partial v}{\partial t}\in \text{null}\; (M_{\alpha\beta}),
$$
so the null space of $M_{\alpha\beta}$ is invariant in time.
%%$$\eqno \#$$ \vskip 0.3cm
\end{proof}

We now apply Hamilton's strong maximum principle to the evolution
equation of the curvature operator $M_{\alpha\beta}$. Recall
$$
\frac{\partial M_{\alpha\beta}}{\partial t} =\Delta
M_{\alpha\beta}+M_{\alpha\beta}^2+M_{\alpha\beta}^{\#}
$$
where $M_{\alpha\beta}^{\#}
=C_\alpha^{\xi\gamma}C_\beta^{\eta\theta}M_{\xi\eta}M_{\gamma\theta}$.
Suppose we have a solution to the Ricci flow with nonnegative
curvature operator. Then by Theorem 2.2.1, the null space of the
curvature operator $M_{\alpha\beta}$ of the solution has constant
rank and is invariant in time and under parallel translation over
some time interval $0<t<\delta$ . Moreover the null space of
$M_{\alpha\beta}$ must also lie in the null space of
$M_{\alpha\beta}^\#$.

Denote by $(n-k)$ the rank of $M_{\alpha\beta}$ on $0<t<\delta$.
Let us diagonalize $M_{\alpha\beta}$ so that $M_{\alpha\alpha}=0$
if $\alpha\le k$ and $M_{\alpha\alpha}>0$ if $\alpha>k$. Then we
have $M_{\alpha\alpha}^\#=0$ also for $\alpha\le k$ from the
evolution equation of $M_{\alpha\alpha}$. Since
$$
0=M_{\alpha\alpha}^\#
=C_\alpha^{\xi\gamma}C_\alpha^{\eta\theta}M_{\xi\eta}M_{\gamma\theta},
$$
it follows that
\begin{align*}
C_\alpha^{\xi\gamma}&  = \<v^\alpha,[v^\xi,v^\gamma]\>\\
&  = 0, \quad \text{if}\quad \alpha\le k\; \text{ and }\;
\xi,\gamma> k.
\end{align*}
This says that the image of $M_{\alpha\beta}$ is a Lie subalgebra
(in fact it is the subalgebra of the restricted holonomy group by
using the Ambrose-Singer holonomy theorem \cite{As53}). This
proves the following result.

%\vskip 0.2cm \noindent {\bf Theorem 2.2.2} (Hamilton \cite{Ha86})\
%\emph{
% CZ-BASELINE source-line:3527 source-kind:theorem source-id:2.2.2
\begin{theorem}[parallel invariance of the curvature-operator image]
  \label{thm:cz-parallel-invariance-of-the-curvature-operator-image}
  \dcref{cao-zhu:2.2.2}
  \group{Cao--Zhu Strong Maximum Principle}
  \source{cao-zhu}
Suppose the curvature operator $M_{\alpha\beta}$ of the initial
metric is nonnegative. Then, under the Ricci flow, for some
interval $0<t<\delta$ the image of $M_{\alpha\beta}$ is a Lie
subalgebra of $so(n)$ which has constant rank and is invariant
under parallel translation and invariant in time.
%$$\eqno\#$$ \vskip 1cm
\end{theorem}


\section{Advanced Maximum Principle for Tensors}

In this section we present Hamilton's advanced maximum principle
\cite{Ha86} for tensors which generalizes Lemma 2.1.3 and shows
how a tensor evolving by a nonlinear heat equation may be
controlled by a system of ODEs. Our presentation in this section
follows closely the papers of Hamilton \cite{Ha86} and Chow-Lu
\cite{Cl04}. An important application of the advanced maximum
principle is the Hamilton-Ivey curvature pinching estimate for the
Ricci flow on three-manifolds given in the next section. More
applications will be given in Chapter 5.

Let $M$ be a complete manifold equipped with a one-parameter
family of Riemannian metrics $g_{ij}(t)$, $0\le t\le T$, with
$T<+\infty$. Let $V\rightarrow M$ be a vector bundle with a
time-independent bundle metric $h_{ab}$ and $\Gamma(V)$ be the
vector space of $C^\infty$ sections of $V$. Let
$$
\nabla_t:\quad \Gamma(V)\rightarrow\Gamma(V\otimes T^*M), \quad
t\in[0,T]
$$
be a smooth family of time-dependent connections compatible with
$h_{ab}$, i.e.
$$
(\nabla_t)_ih_{ab}\stackrel{\Delta}{=}
(\nabla_t)_{\frac{\partial}{\partial x^i}}h_{ab}=0,
$$
for any local coordinate $\{\frac{\partial}{\partial
x^1},\ldots,\frac{\partial}{\partial x^n}\}.$ The Laplacian
$\Delta_t$ acting on a section $\sigma\in\Gamma(V)$ is defined by
$$
\Delta_t\sigma=g^{ij}(x,t)(\nabla_t)_i(\nabla_t)_j\sigma.
$$
For the application to the Ricci flow, we will always assume that
the metrics $g_{ij}(\cdot,t)$ evolve by the Ricci flow. Since $M$
may be noncompact, we assume that, for the sake of simplicity, the
curvature of $g_{ij}(t)$ is uniformly bounded on $M\times [0,T]$.

Let $N: V\times [0,T]\rightarrow V$ be a fiber preserving map,
i.e., $N(x,\sigma,t)$ is a  time-dependent vector field defined on
the bundle $V$ and tangent to the fibers. We assume that
$N(x,\sigma,t)$ is continuous in $x,t$ and satisfies
$$
|N(x,\sigma_1,t) - N(x,\sigma_2,t)| \leq C_B |\sigma_1 -\sigma_2|
$$
for all $ x \in M$, $t\in [0,T]$ and $|\sigma_1|\leq B,
|\sigma_2|\leq B$, where $C_B$ is a positive constant depending
only on $B$. Then we can form the nonlinear heat equation
\begin{equation*}
\frac{\partial}{\partial t}\sigma(x,t)=
\Delta_t\sigma(x,t)+u^i(\nabla_t)_i\sigma(x,t)+N(x,\sigma(x,t),t)
\tag{PDE}
\end{equation*}
where $u^i=u^i(t)$ is a time-dependent vector field on $M$ which
is uniformly bounded on $M\times[0,T]$. Let $K$ be a closed subset
of $V$. One important question is under what conditions will
solutions of the PDE which start in $K$ remain in $K$. To answer
this question, Hamilton \cite{Ha86} imposed the following two
conditions on $K$:
\begin{itemize}
\item[(H1)] $K$ is invariant under parallel translation defined by
the connection $\nabla_t$ for each $t\in[0,T]$; \item[(H2)] in
each fiber $V_x$, the set $K_x\stackrel{\Delta}{=}V_x\cap K$ is
closed and convex.
\end{itemize}

\noindent Then one can judge the behavior of the PDE by comparing
to that of the following ODE
\begin{equation*}
\frac{d\sigma_x}{dt}=N(x,\sigma_x,t)\tag{ODE}
\end{equation*}
for $\sigma_x=\sigma_x(t)$ in each fiber $V_x$. The following
version of Hamilton's advanced maximum principle is from Chow-Lu
\cite{Cl04}

%%\noindent {\bf Theorem 2.3.1} (\textbf
% CZ-BASELINE source-line:3613 source-kind:theorem source-id:2.3.1
\begin{theorem}[advanced maximum principle for time-dependent tensors]
  \label{thm:cz-advanced-maximum-principle-for-time-dependent-tensors}
  \dcref{cao-zhu:2.3.1}
  \group{Cao--Zhu Advanced Maximum Principle for Tensors}
  \source{cao-zhu}
 Let $K$ be a closed subset of $V$ satisfying the
hypothesis {\rm (H1)} and {\rm (H2)}. Suppose that for any $x\in
M$ and any initial time $t_0\in[0,T)$, any solution $\sigma_x(t)$
of the {\rm (ODE)} which starts in $K_x$ at $t_0$ will remain in
$K_x$ for all later times. Then for any initial time $t_0\in
[0,T)$ the solution $\sigma(x,t)$ of the {\rm (PDE)} will remain
in $K$ for all later times provided $\sigma(x,t)$ starts in $K$ at
time $t_0$ and $\sigma(x,t)$ is uniformly bounded with respect to
the bundle metric $h_{ab}$ on $M\times[t_0,T]$.
\end{theorem}

We remark that Lemma 2.1.3 is a special case of the above theorem
where $V$ is given by a symmetric tensor product of a vector
bundle and $K$ corresponds to the convex set consisting of all
nonnegative symmetric bilinear forms. We also remark that Hamilton
\cite{Ha86} established the above theorem for a general evolving
metric $g_{ij}(x,t)$ which does not necessarily satisfy the Ricci
flow.

Before proving Theorem 2.3.1, we need to establish three lemmas in
\cite{Ha86}. Let $\varphi:[a,b]\rightarrow \mathbb{R}$ be a
Lipschitz function. We consider $\frac{d\varphi}{dt}(t)$ at
$t\in[a,b)$ in the sense of limsup of the forward difference
quotients, i.e.,
$$
\frac{d\varphi}{dt}(t) =\limsup_{h\rightarrow
0^+}\frac{\varphi(t+h)-\varphi(t)}{h}.
$$

%\vskip 0.2cm \noindent{\bf Lemma 2.3.2}\ \emph{ \
% CZ-BASELINE source-line:3645 source-kind:lemma source-id:2.3.2
\begin{lemma}[vanishing from a Lipschitz differential inequality]
  \label{lem:cz-vanishing-from-a-lipschitz-differential-inequality}
  \dcref{cao-zhu:2.3.2}
  \group{Cao--Zhu Advanced Maximum Principle for Tensors}
  \source{cao-zhu}
Suppose $\varphi: [a,b]\rightarrow \mathbb{R}$ is Lipschitz
continuous and suppose for some constant $C<+\infty$,
\begin{align*}
&  \frac{d}{dt}\varphi(t)\le C\varphi(t),\quad whenever\ \
\varphi(t)\ge 0\;\ on\ \;[a,b),\\[2mm]
\text{and}\hskip 1cm &  \varphi(a)\le 0.
\end{align*}
Then $\varphi(t)\le 0$ on $[a,b]$.
\end{lemma}

%\vskip 0.1cm\noindent {\bf Proof.}\
\begin{proof}
By replacing $\varphi$ by $e^{-Ct}\varphi$, we may assume
\begin{align*}
&  \frac{d}{dt}\varphi(t)\le 0,\quad \text{whenever}\quad
\varphi(t)\ge 0\;\mbox{ on }\;[a,b),\\
\mbox{and }\hskip 1cm &  \varphi(a)\le 0.
\end{align*}
For arbitrary $\varepsilon>0$, we shall show $\varphi(t)\le
\varepsilon (t-a)$ on $[a,b]$. Clearly we may assume
$\varphi(a)=0$. Since
$$
\limsup_{h\rightarrow 0^+}\frac{\varphi(a+h)-\varphi(a)}{h}\le 0,
$$
there must be some interval $a\le t<\delta$ on which
$\varphi(t)\le \varepsilon (t-a)$.

Let $a\le t<c$ be the largest interval with $c\le b$ such that
$\varphi(t)\le \varepsilon (t-a)$ on $[a,c)$. Then by continuity
$\varphi(t)\le \varepsilon (t-a)$ on the closed interval $[a,c]$.
We claim that $c=b$. Suppose not, then we can find $\delta>0$ such
that $\varphi(t)\le \varepsilon (t-a)$ on $[a,c+\delta]$ since
$$
\limsup_{h\rightarrow 0^+}\frac{\varphi(c+h)-\varphi(c)}{h}\le 0.
$$
This contradicts the choice of the largest interval $[a,c)$.
Therefore, since $\varepsilon>0$ can be arbitrary small, we have
proved $\varphi(t)\le 0$ on $[a,b]$.
\end{proof}
%$$\eqno \#$$ \vskip 0.3cm

The second lemma below is a general principle on the derivative of
a sup-function which will bridge solutions between ODEs and PDEs.
Let $X$ be a complete smooth manifold and $Y$ be a compact subset
of $X$. Let $\psi(x,t)$ be a smooth function on $X\times[a,b]$ and
let $\varphi (t)=\sup\{\psi(y,t)\ |\  y\in Y\}$. Then it is clear
that $\varphi(t)$ is Lipschitz continuous. We have the following
useful estimate on its derivative.

%\vskip 0.2cm\noindent {\bf Lemma 2.3.3}
% CZ-BASELINE source-line:3696 source-kind:lemma source-id:2.3.3
\begin{lemma}[derivative of a fiberwise supremum]
  \label{lem:cz-derivative-of-a-fiberwise-supremum}
  \dcref{cao-zhu:2.3.3}
  \group{Cao--Zhu Advanced Maximum Principle for Tensors}
  \source{cao-zhu}
$$
\frac{d}{dt}\varphi(t) \le \sup\left\{\frac{\partial
\psi}{\partial t}(y,t)\ |\ y\in Y \; \mbox{ satisfies }\;
\psi(y,t)=\varphi(t)\right\}.
$$
\end{lemma}

%\vskip 0.1cm \noindent{\bf Proof.}\
\begin{proof}
Choose a sequence of times $\{t_j\}$ decreasing to $t$ for which
$$
\lim_{t_j\rightarrow t}\frac{\varphi(t_j)-\varphi(t)}{t_j-t}
=\frac{d\varphi(t)}{dt}.
$$
Since $Y$ is compact, we can choose $y_j\in Y$ with
$\varphi(t_j)=\psi(y_j,t_j)$. By passing to a subsequence, we can
assume $y_j\rightarrow y$ for some $y\in Y$. By continuity, we
have $\varphi(t)=\psi(y,t)$. It follows that $\psi(y_j,t)\le
\psi(y,t)$, and then
\begin{align*}
\varphi(t_j)-\varphi(t)&  \le \psi(y_j,t_j)-\psi(y_j,t)\\[2mm]
&  = \frac{\partial}{\partial t}\psi(y_j,\tilde{t}_j)\cdot(t_j-t)
\end{align*}
for some $\tilde{t}_j\in[t,t_j]$ by the mean value theorem. Thus
we have
$$
\lim_{t_j\rightarrow t}\frac{\varphi(t_j)-\varphi(t)}{t_j-t}
\le\frac{\partial}{\partial t}\psi(y,t).
$$
This proves the result.
\end{proof}
%$$\eqno \#$$ \vskip 0.3cm

We remark that the above two lemmas are somewhat standard facts in
the theory of PDEs and we have implicitly used them in the
previous sections when we apply the maximum principle. The third
lemma gives a characterization of when a system of ODEs preserve
closed convex sets in Euclidean space. We will use the version
given in \cite{Cl04}. Let $Z\subset \mathbb R^n$ be a closed
convex subset. We define the {\bf tangent cone}\index{tangent
cone} $T_\varphi Z$ to the closed convex set $Z$ at a point
$\varphi\in
\partial Z$ as the smallest closed convex cone with vertex at
$\varphi$ which contains $Z$.

%\vskip 0.2cm \noindent {\bf Lemma 2.3.4}\ \emph{
% CZ-BASELINE source-line:3743 source-kind:lemma source-id:2.3.4
\begin{lemma}[invariance criterion for closed convex sets]
  \label{lem:cz-invariance-criterion-for-closed-convex-sets}
  \dcref{cao-zhu:2.3.4}
  \group{Cao--Zhu Advanced Maximum Principle for Tensors}
  \source{cao-zhu}
Let $U\subset R^n$ be an open set and $Z\subset U$ be a closed
convex subset. Consider the {\rm ODE} \be
\frac{d\varphi}{dt}=N(\varphi,t)%\eqno(2.3.1)
\ee where $N:\ U\times[0,T]\rightarrow R^n$ is continuous and
Lipschitz in $\varphi$. Then the following two statements are
equivalent.
\begin{itemize}
\item[(i)] For any initial time $t_0\in[0,T]$, any solution of the
{\rm ODE} $(2.3.1)$ which starts in Z at $t_0$ will remain in Z
for all later times; \item[(ii)] $\varphi+N(\varphi,t)\in
T_\varphi Z$ for all $\varphi\in\partial Z$ and $t\in[0,T)$.
\end{itemize}
\end{lemma}

%\vskip 0.1cm \noindent{\bf Proof.}\
\begin{proof}
  \uses{lem:cz-derivative-of-a-fiberwise-supremum,lem:cz-vanishing-from-a-lipschitz-differential-inequality}
We say that a linear function $l$ on $\mathbb R^n$ is a
\textbf{support function}\index{support function} for $Z$ at
$\varphi\in\partial Z$ and write $l\in S_\varphi Z$ if $|l|=1$ and
$l(\varphi)\ge l(\eta)$ for all $\eta \in Z$. Then
$\varphi+N(\varphi,t)\in T_\varphi Z$ if and only if
$l(N(\varphi,t))\le 0$ for all $l\in S_\varphi Z$. Suppose
$l(N(\varphi,t))> 0$ for some $\varphi\in
\partial Z$ and some $l\in S_\varphi Z.$ Then
$$
\frac{d}{dt}l(\varphi)=l\(\frac{d\varphi}{dt}\)=l(N(\varphi,t))>0,
$$
so $l(\varphi)$ is strictly increasing and the solution
$\varphi(t)$ of the ODE (2.3.1) cannot remain in $Z$.

To see the converse, first note that we may assume $Z$ is compact.
This is because we can modify the vector field $N(\varphi,t)$ by
multiplying a cutoff function which is everywhere nonnegative,
equals one on a large ball and equals zero on the complement of a
larger ball. The paths of solutions of the ODE are unchanged
inside the first large ball, so we can intersect $Z$ with the
second ball to make $Z$ convex and compact. If there were a
counterexample before the modification there would still be one
after as we chose the first ball large enough.

Let $s(\varphi)$ be the distance from $\varphi$ to $Z$ in
$\mathbb{R}^n$. Clearly $s(\varphi)=0$ if $\varphi\in Z$. Then
$$
s(\varphi)=\sup\{l(\varphi-\eta)\ |\ \eta\in\partial Z\; \text{
and }\; l\in S_\eta Z\}.
$$
The sup is taken over a compact subset of $\mathbb R^n\times
\mathbb R^n$. Hence by Lemma 2.3.3
$$
\frac{d}{dt}s(\varphi)\le \sup\{l(N(\varphi,t))\ |\
\eta\in\partial Z, l\in S_\eta Z\; \text{ and }\;
s(\varphi)=l(\varphi-\eta)\}.
$$
It is clear that the sup on the RHS of the above inequality can be
takeen only when $\eta$ is the unique closest point in $Z$ to
$\varphi$ and $l$ is the linear function of length one with
gradient in the direction of $\varphi-\eta$. Since $N(\varphi,t)$
is Lipschitz in $\varphi$ and continuous in $t$, we have
$$
|N(\varphi,t)-N(\eta,t)|\le C|\varphi-\eta|
$$
for some constant $C$ and all $\varphi$ and $\eta$ in the compact
set $Z$.

By hypothesis (ii), $$l(N(\eta,t))\le 0,$$ and for the unique
$\eta$, the closest point in $Z$ to $\varphi$,
$$
|\varphi-\eta|=s(\varphi).
$$
Thus
\begin{align*}
\frac{d}{dt}s(\varphi) &  \le \sup \left\{\begin{array}{c}
l(N(\eta,t))+ |l(N(\varphi,t))-l(N(\eta,t))| \quad |\quad
\eta\in\partial Z,\vspace*{3mm}\\
 l\in S_\eta Z,\quad \text{and}\quad
s(\varphi)=l(\varphi-\eta)
\end{array}\right\} \\
&  \le Cs(\varphi).
\end{align*}
Since $s(\varphi)=0$ to start at $t_0$, it follows from Lemma
2.3.2 that $s(\varphi)=0$ for $t\in[t_0,T]$. This proves the
lemma.
\end{proof}
%$$\eqno \#$$ \vskip 0.3cm

We are now ready to prove Theorem 2.3.1.

%\vskip 0.1cm \noindent
\medskip
{\bf\em Proof of Theorem} {\bf 2.3.1.} \ Since the solution
$\sigma(x,t)$ of the (PDE) is uniformly bounded with respect to
the bundle metric $h_{ab}$ on $M\times [t_0,T]$ by hypothesis, we
may assume that $K$ is contained in a tubular neighborhood $V(r)$
of the zero section in $V$ whose intersection with each fiber
$V_x$ is a ball of radius $r$ around the origin measured by the
bundle metric $h_{ab}$ for some large $r>0$.

Recall that $g_{ij}(\cdot,t), t\in[0,T]$, is a smooth solution to
the Ricci flow with uniformly bounded curvature on $M\times
[0,T]$. From Lemma 2.1.1, we have a smooth function $f$ such that
$f\geq 1$ everywhere, $f(x)\rightarrow +\infty$ as
$d_0(x,x_0)\rightarrow +\infty$ for some fixed point $x_0\in M$,
and the first and second covariant derivatives with respect to the
metrics $g_{ij}(\cdot,t)$ are uniformly bounded on $M\times
[0,T]$. Using the metric $h_{ab}$ in each fiber $V_x$ and writing
$|\varphi- \eta|$ for the distance between $\varphi \in V_x$ and
$\eta \in V_x$, we set
$$
s(t)=\sup_{x\in M}\{ \inf\{|\sigma(x,t)-\eta|\  | \ \eta\in
K_x\stackrel{\Delta}{=} K\cap V_x\}-\epsilon e^{At}f(x)\}
$$
where $\epsilon$ is an arbitrarily small positive number and $A$
is a positive constant to be determined. We rewrite the function
$s(t)$ as
$$
s(t)=\sup\{l(\sigma(x,t)-\eta)-\epsilon e^{At}f(x)\ |\  x\in M,
\eta\in \partial K_x  \mbox{ and }  l\in S_\eta K_x\}.
$$
By the construction of the function $f$, we see that the sup is
taken in a compact subset of $M\times V \times V^*$ for all $t$.
Then by Lemma 2.3.3, \be \frac{ds(t)}{dt}\leq \sup\left\{
\frac{\partial}{\partial
t}l(\sigma(x,t)-\eta)-\epsilon Ae^{At}f(x)\right\} %\eqno (2.3.2)
\ee where the sup is over all $x\in M, \eta\in\partial K_x$ and $
l \in S_\eta K_x$ such that
$$
l(\sigma(x,t)-\eta)-\epsilon e^{At}f(x)=s(t);
$$
in particular we have $|\sigma(x,t)-\eta|=l(\sigma(x,t)-\eta)$,
where $\eta$ is the unique closest point in $K_x$ to
$\sigma(x,t)$, and $l$ is the linear function of length one on the
fiber $V_x$ with gradient in the direction of $\eta$ to
$\sigma(x,t)$. We compute at these $(x,\eta,l)$,
\begin{align}
& \frac{\partial}{\partial t}l(\sigma(x,t)-\eta)-\epsilon A
e^{At}f(x)\\
&  =l\(\frac{\partial \sigma(x,t)}{\partial t}\)
-\epsilon A e^{At}f(x)\nn\\
&  =l(\Delta_t \sigma(x,t))
+l(u^i(x,t)(\nabla_t)_i\sigma(x,t))+l(N(x,\sigma(x,t),t))-\epsilon
A e^{At}f(x).\nn
\end{align} %\eqno(2.3.3)
By the assumption and Lemma 2.3.4 we have $\eta+N(x,\eta,t)\in
T_\eta K_x$. Hence, for those $(x,\eta,l)$, $l(N(x,\eta,t))\leq 0$
and then
\begin{align}
& l(N(x,\sigma(x,t),t))\\
&  \leq l(N(x,\eta,t))+|N(x,\sigma(x,t),t)-N(x,\eta,t)| \nn\\
&  \leq C|\sigma(x,t)-\eta|=C(s(t)+\epsilon e^{At}f(x)) \nn
\end{align} %\eqno(2.3.4)
for some positive constant $C$ by the assumption that
$N(x,\sigma,t)$ is Lipschitz in $\sigma$ and the fact that the sup
is taken on a compact set. Thus the combination of
(2.3.2)--(2.3.4) gives \be \frac{ds(t)}{dt}\leq l(\Delta_t
\sigma(x,t)) +l(u^i(x,t)(\nabla_t)_i\sigma(x,t))+Cs(t)+\epsilon
(C-A)e^{At}f(x)
%\eqno(2.3.5)
\ee for those $x\in M, \eta \in \partial K_x$ and $l\in  S_\eta
K_x$ such that $ l(\sigma(x,t)-\eta)-\epsilon e^{At}f(x)=s(t)$.

Next we estimate the first two terms of (2.3.5). As we extend a
vector in a bundle from a point $x$ by parallel translation along
geodesics emanating radially out of $x$, we will get a smooth
section of the bundle in some small neighborhood of $x$ such that
all the symmetrized covariant derivatives at $x$ are zero. Now let
us extend $\eta\in V_x$ and $l\in V^*_x$ in this manner.  Clearly,
we continue to have $|l|(\cdot)=1$. Since $K$ is invariant under
parallel translations, we continue to have $\eta (\cdot)\in
\partial K$ and $l(\cdot)$ as a support function for $K$ at $\eta
(\cdot)$. Therefore
$$
l(\sigma(\cdot,t)-\eta(\cdot))-\epsilon e^{At}f(\cdot)\leq s(t)
$$
in the neighborhood. It follows that the function
$l(\sigma(\cdot,t)-\eta(\cdot))-\epsilon e^{At}f(\cdot)$ has a
local maximum at $x$, so at $x$
\begin{align*}
(\nabla _t)_i(l(\sigma(x,t)-\eta)-\epsilon e^{At}f(x))&=0,\\
\text{and }\; \Delta_t(l(\sigma(x,t)-\eta)-\epsilon e^{At}f(x))&
\leq 0.
\end{align*}
Hence at $x$
\begin{align*}
l((\nabla _t)_i\sigma(x,t))-\epsilon e^{At}(\nabla_t)_if(x)&=0,\\
\text{and }\; l(\Delta_t \sigma(x,t))-\epsilon e^{At}\Delta_t
f(x)&\leq 0.
\end{align*}
Therefore by combining with (2.3.5), we have
\begin{align*}
\frac{d}{dt}s(t)&  \leq Cs(t)+\epsilon
(\Delta_tf(x)+u^i(\nabla_t)_if(x)+(C-A)f(x))e^{At} \\
&  \leq Cs(t)
\end{align*}
for $A>0$ large enough, since $f(x)\geq 1$ and the first and
second covariant derivatives of $f$ are uniformly bounded on
$M\times[0,T]$. So by applying Lemma 2.3.2 and the arbitrariness
of $\epsilon$, we have completed the proof of Theorem 2.3.1.
\endproof
%$$\eqno \#$$ \vskip 0.3cm

Finally, we would like to state a useful generalization of Theorem
2.3.1 by Chow and Lu in \cite{Cl04} which allows the set $K$ to
depend on time. One can consult the paper \cite{Cl04} for the
proof.

%\vskip 0.2cm\noindent {\bf Theorem 2.3.5} (Chow and Lu
%\cite{Cl04})\ \emph{
% CZ-BASELINE source-line:3951 source-kind:theorem source-id:2.3.5
\begin{theorem}[maximum principle for closed fiberwise convex sets]
  \label{thm:cz-maximum-principle-for-closed-fiberwise-convex-sets}
  \dcref{cao-zhu:2.3.5}
  \group{Cao--Zhu Advanced Maximum Principle for Tensors}
  \source{cao-zhu}
Let $K(t)\subset V$, $t\in [0,T]$ be closed subsets which satisfy
the following hypotheses
\begin{itemize}
\item[(H3)] $K(t)$ is invariant under parallel translation defined
by the connection $\nabla_t$ for each $t\in [0,T]$; \item[(H4)] in
each fiber $V_x$, the set $K_x(t)\stackrel{\Delta}{=}K(t)\cap V_x$
is nonempty, closed and convex for each $t\in [0,T]$; \item[(H5)]
the space-time track $\bigcup\limits_{t\in[0,T]}(\partial
K(t)\times\{t\})$ is a closed subset of $V\times[0,T]$.
\end{itemize}
Suppose that, for any $x\in M$ and any initial time $t_0\in[0,T)$,
and for any solution $\sigma_x(t)$ of the {\rm (ODE)} which starts
in $K_x(t_0)$, the solution $\sigma_x(t)$ will remain in $K_x(t)$
for all later times. Then for any initial time $t_0\in[0,T)$ the
solution $\sigma(x,t)$ of the {\rm (PDE)} will remain in $K(t)$
for all later times if $\sigma(x,t)$ starts in $K(t_0)$ at time
$t_0$ and the solution $\sigma(x,t)$ is uniformly bounded with
respect to the bundle metric $h_{ab}$ on $M\times[t_0,T]$.
\end{theorem}
%$$\eqno \#$$ \vskip 1cm


\section{Hamilton-Ivey Curvature Pinching Estimate}

The Hamilton-Ivey curvature pinching estimate \cite{Ha95F, Iv}
roughly says that if a solution to the Ricci flow on a
three-manifold becomes singular (i.e., the curvature goes to
infinity) as time $t$ approaches the maximal time $T$, then the
most negative sectional curvature will be small compared to the
most positive sectional curvature. This pinching estimate plays a
crucial role in analyzing the formation of singularities in the
Ricci flow on three-manifolds. The proof here is based on the
argument in Hamilton \cite{Ha95F}. The estimate was later improved
by Hamilton \cite{Ha99} which will be presented in Section 5.3
(see Theorem 5.3.2).

Consider a complete solution to the Ricci flow
$$
\frac{\partial}{\partial t}g_{ij}=-2R_{ij}
$$
on a complete three-manifold with bounded curvature in space for
each time $t\geq 0$. Recall from Section 1.3 that the evolution
equation of the curvature operator $M_{\alpha\beta}$ is given by
\be \frac{\partial}{\partial t}M_{\alpha\beta}=\Delta
M_{\alpha\beta}+M_{\alpha\beta}^2+M_{\alpha\beta}^\# %\eqno(2.4.1)
\ee where $M_{\alpha\beta}^2$ is the operator square
$$
M_{\alpha\beta}^2=M_{\alpha\gamma}M_{\beta\gamma}
$$
and $M_{\alpha\beta}^\#$ is the Lie algebra $so(n)$ square
$$
M_{\alpha\beta}^\#=C_\alpha^{\gamma\zeta}
C_\beta^{\eta\theta}M_{\gamma\eta}M_{\zeta\theta}.
$$
In dimension $n=3$, we know that $M_{\alpha\beta}^\#$ is the
adjoint matrix of $M_{\alpha\beta}$. If we diagonalize
$M_{\alpha\beta}$ with eigenvalues $\lambda \ge \mu \ge \nu$ so
that
$$(M_{\alpha\beta})=\left(
\begin{array}{ccc}\lambda&  \ &  \ \\
\ &  \mu &  \ \\
\ &  \ &  \nu
\end{array}
\right),
$$
then $M_{\alpha\beta}^2$ and $M_{\alpha\beta}^\#$ are also
diagonal, with
$$
(M_{\alpha\beta}^2)=\left(
\begin{array}{ccc}\lambda^2&  \ &  \ \\
\ &  \mu^2 &  \ \\
\ &  \ &  \nu^2
\end{array}
\right) \ \ \text{and} \ \ (M_{\alpha\beta}^\#)=\left(
\begin{array}{ccc}\mu\nu&  \ &  \ \\
\ &  \lambda\nu &  \ \\
\ &  \ &  \lambda\mu
\end{array}
\right).
$$

Thus the ODE corresponding to PDE (2.4.1) for $M_{\alpha\beta}$
(in the space of $3\times 3 $ matrices) is given by the following
system \be
      \left\{
       \begin{array}{lll}
  \frac{d}{dt}\lambda=\lambda^2+\mu\nu,\\[4mm]
  \frac{d}{dt}\mu=\mu^2+\lambda\nu,\\[4mm]
  \frac{d}{dt}\nu=\nu^2+\lambda\mu.
       \end{array}
    \right.
%  \eqno(2.4.2)
\ee

Let $P$ be the principal bundle of the manifold and form the
associated bundle $V=P\times_GE$, where $G=O(3)$ and $E$ is the
vector space of symmetric bilinear forms on $so(3)$. The curvature
operator $M_{\alpha\beta}$ is a smooth section of $V=P\times_GE$.
According to Theorem 2.3.1, any closed convex set of curvature
operator matrices $M_{\alpha\beta}$ which is $O(3)$-invariant (and
hence invariant under parallel translation) and preserved by ODE
(2.4.2) is also preserved by the Ricci flow.

We are now ready to state and prove the \textbf{Hamilton-Ivey
pinching estimate} \index{Hamilton-Ivey pinching estimate}.

%\vskip 0.2cm \noindent{\bf Theorem 2.4.1} (Hamilton \cite{Ha95F},
%Ivey \cite{Iv})\ \emph{
% CZ-BASELINE source-line:4060 source-kind:theorem source-id:2.4.1
\begin{theorem}[Hamilton–Ivey curvature pinching]
  \label{thm:cz-hamiltonivey-curvature-pinching}
  \dcref{cao-zhu:2.4.1}
  \group{Cao--Zhu Hamilton-Ivey Curvature Pinching Estimate}
  \source{cao-zhu}
Suppose we have a solution to the Ricci flow on a three-manifold
which is complete with bounded curvature for each $t\geq0$. Assume
at $t=0$ the eigenvalues $\lambda\geq\mu\geq\nu$ of the curvature
operator at each point are bounded below by $\nu\geq-1$. The
scalar curvature $R=\lambda+\mu+\nu$ is their sum. Then at all
points and all times $t\geq0$ we have the pinching estimate
$$
R\geq (-\nu)[\log(-\nu)-3],
$$
whenever $\nu<0$.
\end{theorem}

%\vskip 0.1cm \noindent{\bf Proof. }
\begin{proof} The proof is taken from Hamilton \cite{Ha95F}.
  \uses{thm:cz-advanced-maximum-principle-for-time-dependent-tensors}
Consider the function
$$
y=f(x)=x(\log x-3)
$$
defined on $e^2\leq x<+\infty$. It is easy to check that $f$ is
increasing and convex  with range $-e^2\leq y<+\infty$. Let
$f^{-1}(y)=x$ be the inverse function, which is also increasing
but concave and satisfies \be
\lim_{y\rightarrow\infty}\frac{f^{-1}(y)}{y}=0 %\eqno(2.4.3)
\ee Consider also the set $K$ of matrices $M_{\alpha\beta}$
defined by the inequalities \be K:
   \left  \{
     \begin{array}{lll}
         \lambda+\mu+\nu\geq-3,\\[4mm]
         \nu+f^{-1}(\lambda+\mu+\nu)\geq0.
     \end{array}
     \right. %\eqno(2.4.4)
\ee By Theorem 2.3.1 and the assumptions in Theorem 2.4.1 at
$t=0$, we only need to check that the set $K$ defined above is
closed, convex and preserved by the ODE (2.4.2).

Clearly $K$ is closed because $f^{-1}$ is continuous.
$\lambda+\mu+\nu$ is just the trace function of $3\times 3$
matrices which is a linear function. Hence the first inequality in
(2.4.4) defines a linear half-space, which is convex. The function
$\nu$ is the least eigenvalue function, which is concave. Also
note that $f^{-1}$ is concave. Thus the second inequality in
(2.4.4) defines a convex set as well. Therefore we proved $K$ is
closed and convex.

Under the ODE (2.4.2)
\begin{align*}
    \frac{d}{dt}(\lambda+\mu+\nu)
&  = \lambda^2+\mu^2+\nu^2+\lambda\mu+\lambda\nu+\mu\nu\\
&  = \frac{1}{2}[(\lambda+\mu)^2+(\lambda+\nu)^2+(\mu+\nu)^2]\\
&  \geq 0.
\end{align*}
Thus the first inequality in (2.4.4) is preserved by the ODE.

The second inequality in (2.4.4) can be written as
$$
\lambda+\mu+\nu\geq f(-\nu),\ \ \ \text{whenever}\ \nu\leq-e^2,
$$
which becomes \be \lambda+\mu\geq (-\nu)[\log(-\nu)-2],\ \ \
\text{whenever}\
\nu\leq-e^2.%\eqno(2.4.5)
\ee To show the inequality is preserved we only need to look at
points on the boundary of the set. If
$\nu+f^{-1}(\lambda+\mu+\nu)=0$ then
$\nu=-f^{-1}(\lambda+\mu+\nu)\le -e^2$ since $f^{-1}(y)\ge e^2$.
Hence the RHS of (2.4.5) is nonnegative. We thus have
$\lambda\geq0$ because $\lambda\geq\mu$. But $\mu$ may have either
sign. We split our consideration into two cases:

\medskip
{\it Case} (i): $\mu\geq0$.

\smallskip
We need to verify
$$
\frac{d\lambda}{dt}+\frac{d\mu}{dt}\geq(\log(-\nu)-1)\frac{d(-\nu)}{dt}
$$
when $\lambda+\mu=(-\nu)[\log(-\nu)-2]$. Solving for
$$
\log(-\nu)-2=\frac{\lambda+\mu}{(-\nu)}
$$
and substituting above, we must show
$$
\lambda^2+\mu\nu+\mu^2+\lambda\nu
\geq\(\frac{\lambda+\mu}{(-\nu)}+1\)(-\nu^2-\lambda\mu)
$$
which is equivalent to
$$
(\lambda^2+\mu^2)(-\nu)+\lambda\mu(\lambda+\mu+(-\nu))+(-\nu)^3\geq
0.
$$
Since $\lambda,\mu$ and $(-\nu)$ are all nonnegative we are done
in the first case.

\medskip
{\it Case} (ii): $\mu<0$.

\smallskip
We need to verify
$$
\frac{d\lambda}{dt}\geq\frac{d(-\mu)}{dt}
+(\log(-\nu)-1)\frac{d(-\nu)}{dt}
$$
when $\lambda=(-\mu)+(-\nu)[\log(-\nu)-2]$. Solving for
$$
\log(-\nu)-2=\frac{\lambda-(-\mu)}{(-\nu)}
$$
and substituting above, we need to show
$$
\lambda^2+\mu\nu\geq-\mu^2-\lambda\nu
+\(\frac{\lambda-(-\mu)}{(-\nu)}+1\)(-\nu^2-\lambda\mu)
$$
or
$$
\lambda^2+(-\mu)(-\nu)\geq\lambda(-\nu)-(-\mu)^2
+\(\frac{\lambda-(-\mu)}{(-\nu)}+1\)(\lambda(-\mu)-(-\nu)^2)
$$
which reduces to
$$
\lambda^2(-\nu)+\lambda(-\mu)^2+(-\mu)^2(-\nu)+(-\nu)^3
\geq\lambda^2(-\mu)+\lambda(-\mu)(-\nu)
$$
or equivalently
$$
(\lambda^2-\lambda(-\mu)+(-\mu)^2)((-\nu)-(-\mu))
+(-\mu)^3+(-\nu)^3\geq0.
$$
Since $\lambda^2-\lambda(-\mu)+(-\mu)^2\geq0\ \ \text{and}\ \
(-\nu)-(-\mu)\geq0$ we are also done in the second case.

Therefore the proof is completed.
\end{proof}
%$$\eqno \#$$ \vskip 1cm


\section{Li-Yau-Hamilton Estimates}

In \cite{LY}, Li-Yau developed a fundamental gradient estimate,
now called Li-Yau estimate, for positive solutions to the heat
equation on a complete Riemannian manifold with nonnegative Ricci
curvature. They used it to derive the Harnack inequality for such
solutions by path integration. Then based on the suggestion of
Yau, Hamilton \cite{Ha88} developed a similar estimate for the
scalar curvature of solutions to the Ricci flow on a Riemann
surface with positive curvature, and later obtained a matrix
version of the Li-Yau estimate \cite{Ha93} for solutions to the
Ricci flow with positive curvature operator in all dimensions.
This matrix version of the Li-Yau estimate is the
\textbf{Li-Yau-Hamilton estimate}\index{Li-Yau-Hamilton estimate},
which we will present in this section. Most of the presentation
follows the original papers of Hamilton \cite{Ha88, Ha93, Ha93E}.

We have seen that in the Ricci flow the curvature tensor satisfies
a nonlinear heat equation, and the nonnegativity of the curvature
operator is preserved by the Ricci flow. Roughly speaking, the
Li-Yau-Hamilton estimate says the nonnegativity of a certain
combination of the derivatives of the curvature up to second order
is also preserved by the Ricci flow. This estimate plays a central
role in the analysis of formation of singularities and the
application of the Ricci flow to three-manifold topology.

Let us begin by describing the Li-Yau estimate \cite{LY} for
positive solutions to the heat equation on a complete Riemannian
manifold with nonnegative Ricci curvature.

%\vskip 0.2cm \noindent{\bf Theorem 2.5.1} (Li-Yau \cite{LY})\emph{
% CZ-BASELINE source-line:4226 source-kind:theorem source-id:2.5.1
\begin{theorem}[Li–Yau differential Harnack estimate]
  \label{thm:cz-liyau-differential-harnack-estimate}
  \dcref{cao-zhu:2.5.1}
  \group{Cao--Zhu Li-Yau-Hamilton Estimates}
  \source{cao-zhu}
Let $(M, g_{ij})$ be an $n$-dimensional complete Riemannian
manifold with nonnegative Ricci curvature. Let $u(x,t)$ be any
positive solution to the heat equation
$$
\frac{\partial u}{\partial t}=\Delta u \qquad {\text{on}} \ \
M\times [0,\infty).
$$
Then we have \be \frac{\partial u}{\partial t}-\frac{|\nabla
u|^2}{u}+\frac{n}{2t}u\ge 0 \qquad \text{on } \ \ M\times
(0,\infty).%\eqno(2.5.1)
\ee
%$$\eqno \#$$ \vskip 0.3cm
\end{theorem}

We remark that, as observed by Hamilton (cf. \cite{Ha93}), one can
in fact prove that for any vector field $V^i$ on $M$, \be
\frac{\partial u}{\partial t}+2\nabla u\cdot V+u|V|^2
+\frac{n}{2t}u\ge 0. %\eqno(2.5.2)
\ee If we take the optimal vector field $V=-\nabla u/u$, we
recover the inequality (2.5.1).

Now we consider the Ricci flow on a Riemann surface. Since in
dimension two the Ricci curvature is given by
$$
R_{ij}=\frac{1}{2} Rg_{ij},
$$
the Ricci flow (1.1.5) becomes \be
\frac{\partial g_{ij}}{\partial t}=-R g_{ij}.%\eqno(2.5.3)
\ee

Now let $g_{ij}(x,t)$ be a complete solution of the Ricci flow
(2.5.3) on a Riemann surface $M$ and $0\le t<T$. Then the scalar
curvature $ R(x,t)$ evolves by the semilinear equation
$$
\frac{\partial R}{\partial t}=\triangle R+R^{2}
$$
on $M\times[0,T)$. Suppose the scalar curvature of the initial
metric is bounded, nonnegative everywhere and positive somewhere.
Then it follows from Proposition 2.1.2 that the scalar curvature
$R(x,t)$ of the evolving metric remains nonnegative. Moreover,
from the standard strong maximum principle (which works in each
local coordinate neighborhood), the scalar curvature is positive
everywhere for $t>0$. In \cite{Ha88}, Hamilton obtained the
following Li-Yau estimate for the scalar curvature $R(x,t)$.

%\vskip 0.2cm\noindent{\bf Theorem 2.5.2} (Hamilton \cite{Ha88})\
% CZ-BASELINE source-line:4273 source-kind:theorem source-id:2.5.2
\begin{theorem}[surface Ricci-flow differential Harnack estimate]
  \label{thm:cz-surface-ricci-flow-differential-harnack-estimate}
  \dcref{cao-zhu:2.5.2}
  \group{Cao--Zhu Li-Yau-Hamilton Estimates}
  \source{cao-zhu}
Let $g_{ij}(x,t)$ be a complete solution of the Ricci flow on a
surface $M$. Assume the scalar curvature of the initial metric is
bounded, nonnegative everywhere and positive somewhere. Then the
scalar curvature $R(x,t)$ satisfies the Li-Yau estimate \be
\frac{\partial R}{\partial t} - \frac{|\nabla R|^2}{R} +
\frac{R}{t} \geq 0.%\eqno(2.5.4)
\ee
\end{theorem}

%\vskip 0.1cm\noindent{\bf Proof} \
\begin{proof}
  \uses{lem:cz-tensor-maximum-principle-on-complete-manifolds}
By the above discussion, we know $R(x,t)>0$ for $t>0$. If we set
$$
L=\log R(x,t)\quad \text{for}\quad t>0,
$$
then
\begin{align*}
\frac{\partial}{\partial t}L&  = \frac{1}{R}(\triangle
R+R^{2})\\
&  = \triangle L+|\nabla L|^{2}+R
\end{align*}
and (2.5.4) is equivalent to
$$
\frac{\partial L}{\partial t}-|\nabla L|^{2}+\frac{1}{t}
=\triangle L +R +\frac{1}{t}\geq0.
$$

Following Li-Yau \cite{LY} in the linear heat equation case, we
consider the quantity \be Q =\frac{\partial L}{\partial t}-|\nabla
L|^{2}=\triangle L+R.
%\eqno(2.5.5)
\ee Then by a direct computation,
\begin{align*}
\frac{\partial Q}{\partial t}
&  = \frac{\partial}{\partial t}(\triangle L+R)\\
&  = \triangle\(\frac{\partial L}{\partial t}\)+R\triangle
L+\frac{\partial R}{\partial t}\\
&  = \triangle Q+2\nabla L\cdot\nabla Q+2|\nabla ^{2}L|^{2}
+2R(\triangle L)+R^{2}\\
&  \geq \triangle Q+2\nabla L\cdot\nabla Q+Q^{2}.
\end{align*}
So we get
$$
\frac{\partial}{\partial t}
\(Q+\frac{1}{t}\)\geq\triangle\(Q+\frac{1}{t}\)+2\nabla
L\cdot\nabla \(Q+\frac{1}{t}\)+\(Q-\frac{1}{t}\)\(Q+\frac{1}{t}\).
$$
Hence by a similar maximum principle argument as in the proof of
Lemma 2.1.3, we obtain
$$
Q+\frac{1}{t}\ge 0.
$$
This proves the theorem.
%$$\eqno \#$$ \vskip 0.3cm
\end{proof}

As an immediate consequence, we obtain the following Harnack
inequality for the scalar curvature $R$ by taking the Li-Yau type
path integral as in \cite{LY}.

% \vskip 0.2cm\noindent{\bf Corollary 2.5.3} (Hamilton \cite{Ha88})\
% CZ-BASELINE source-line:4335 source-kind:corollary source-id:2.5.3
\begin{corollary}[surface Ricci-flow integrated Harnack estimate]
  \label{cor:cz-surface-ricci-flow-integrated-harnack-estimate}
  \dcref{cao-zhu:2.5.3}
  \group{Cao--Zhu Li-Yau-Hamilton Estimates}
  \source{cao-zhu}
Let $g_{ij}(x,t)$ be a complete solution of the Ricci flow on a
surface with bounded and nonnegative scalar curvature. Then for
any points $x_{1},x_{2}\in M$, and $0<t_{1}<t_{2}$, we have
$$
R(x_2,t_2)
\geq\frac{t_1}{t_2}e^{-d_{t_1}{(x_1,x_2)}^2/{4(t_2-t_1)}}R(x_1,t_1).
$$
\end{corollary}

%\vskip 0.1cm \noindent {\bf Proof.}
\begin{proof}
  \uses{thm:cz-surface-ricci-flow-differential-harnack-estimate}
Take the geodesic path $\gamma(\tau)$, $\tau\in[t_1,t_2],$ from
$x_1$ to $ x_2$ at time $t_1$ with constant velocity
$d_{t_1}(x_1,x_2)/(t_2-t_1).$ Consider the space-time path
$\eta(\tau)=(\gamma(\tau),\tau)$, $\tau\in[t_1,t_2]$. We compute
\begin{align*}
\log\frac{R(x_2,t_2)}{R(x_1,t_1)}
&  = \int^{t_2}_{t_1}\frac{d}{d\tau}L(\gamma(\tau),\tau)d\tau\\
&  = \int^{t_2}_{t_1}\frac{1}{R}\(\frac{\partial R}{\partial \tau}
+\nabla R \cdot \frac{d\gamma}{d\tau}\)d\tau\\
&  \geq \int^{t_2}_{t_1}\(\frac{\partial L}{\partial\tau} -|\nabla
L|^2_{g_{ij}(\tau)}
-\frac{1}{4}\left|\frac{d\gamma}{d\tau}\right|^2_{g_{ij}(\tau)}\)d\tau.
\end{align*}
Then by Theorem 2.5.2 and the fact that the metric is shrinking
(since the scalar curvature is nonnegative), we have
\begin{align*}
\log\frac{R(x_2,t_2)}{R(x_1,t_1)} &
\geq\int^{t_2}_{t_1}\(-\frac{1}{\tau}-
\frac{1}{4}\left|\frac{d\gamma}{d\tau}\right|^2_{g_{ij}(\tau)}\)d\tau \\
&  = \log\frac{t_1}{t_2}-\frac{d_{t_1}(x_1,x_2)^2}{4(t_2-t_1)}
\end{align*}
After exponentiating above, we obtain the desired Harnack
inequality.
\end{proof}
%$$\eqno \#$$ \vskip 0.3cm

To prove a similar inequality as (2.5.4) for the scalar curvature
of solutions to the Ricci flow in higher dimensions is not so
simple. First of all, we will need to require nonnegativity of the
curvature operator (which we know is preserved under the Ricci
flow).  Secondly, one does not get inequality (2.5.4) directly,
but rather indirectly as the trace of certain matrix estimate. The
key ingredient in formulating this matrix version is to derive
some identities from the soliton solutions and prove an elliptic
inequality based on these quantities. Hamilton found such a
general principle which was based on the idea of Li-Yau \cite{LY}
when an identity is checked on the heat kernel before an
inequality was found. To illustrate this point, let us first
examine the heat equation case. Consider the heat kernel
$$
u(x,t)=(4\pi t)^{-n/2} e^{-|x|^2/4t}
$$
for the standard heat equation on $\mathbb R^n$ which can be
considered as an expanding soliton solution.

Differentiating the function $u$, we get \be \nabla_j u=-u
\frac{x_j}{2t} \; \text{ or }\;
\nabla_j u+uV_j=0,%\eqno (2.5.6)
\ee where
$$
V_j=\frac{x_j}{2t}=-\frac{\nabla_j u}{u}.
$$

Differentiating (2.5.6), we have \be \nabla_i\nabla_j
u+\nabla_iuV_j+\frac{u}{2t}\delta_{ij}=0.
%\eqno(2.5.7)
\ee To make the expression in (2.5.7) symmetric in $i, j$, we
multiply $V_i$ to (2.5.6) and add to (2.5.7) and obtain \be
\nabla_i\nabla_j u+\nabla_iuV_j+\nabla_juV_i+uV_iV_j
+\frac{u}{2t}\delta_{ij}=0. %\eqno (2.5.8)
\ee Taking the trace in (2.5.8) and using the equation ${\partial
u}/{\partial t}=\Delta u$, we arrive at
$$
\frac{\partial u}{\partial t} +2\nabla u\cdot V+u|V|^2
+\frac{n}{2t} u=0,
$$
which shows that the Li-Yau inequality (2.5.1) becomes an equality
on our expanding soliton solution $u$! Moreover, we even have the
matrix identity (2.5.8).

Based on the above observation and using a similar process,
Hamilton found a matrix quantity, which vanishes on expanding
gradient Ricci solitons and is nonnegative for any solution to the
Ricci flow with nonnegative curvature operator. Now we describe
the process of finding the Li-Yau-Hamilton quadratic for the Ricci
flow in arbitrary dimension.

Consider a homothetically expanding gradient soliton $g$, we have
\be
R_{ab}+\frac{1}{2t}g_{ab}=\nabla_aV_b %\eqno(2.5.9)
\ee in the orthonormal frame coordinate chosen as in Section 1.3.
Here $V_b=\nabla_bf$ for some function $f$. Differentiating
(2.5.9) and commuting give the first order relations
\begin{align}
\nabla_aR_{bc}-\nabla_bR_{ac}
&  = \nabla_a\nabla_bV_c-\nabla_b\nabla_aV_c\\
&  = R_{abcd}V_d,\nn
\end{align}   %%\eqno(2.5.10)
and differentiating again, we get
\begin{align*}
\nabla_a\nabla_bR_{cd}-\nabla_a\nabla_cR_{bd}
&  = \nabla_a(R_{bcde}V_e)\\
&  = \nabla_aR_{bcde}V_e+R_{bcde}\nabla_aV_e\\
&  = \nabla_aR_{bcde}V_e+R_{ae}R_{bcde}+\frac{1}{2t}R_{bcda}.
\end{align*}
We further take the trace of this on $a$ and $b$ to get
$$
\Delta R_{cd}-\nabla_a\nabla_cR_{ad}-R_{ae}R_{acde}
+\frac{1}{2t}R_{cd}-\nabla_aR_{acde}V_e=0,
$$
and then by commuting the derivatives and second Bianchi identity,
$$
\Delta R_{cd}-\frac{1}{2}\nabla_c\nabla_dR
+2R_{cade}R_{ae}-R_{ce}R_{de}+\frac{1}{2t}R_{cd}+(\nabla
_eR_{cd}-\nabla_dR_{ce})V_e=0.
$$
Let us define
\begin{align*}
M_{ab}&  = \Delta R_{ab}-\frac{1}{2}\nabla_a\nabla_bR
+2R_{acbd}R_{cd}-R_{ac}R_{bc}+\frac{1}{2t}R_{ab},\\
P_{abc}&  = \nabla_aR_{bc}-\nabla_bR_{ac}.
\end{align*}
Then \be
M_{ab}+P_{cba}V_c=0,%\eqno(2.5.11)
\ee We rewrite (2.5.10) as
$$
P_{abc}=R_{abcd}V_d
$$
and then \be
P_{cab}V_c+R_{acbd}V_cV_d=0. %\eqno(2.5.12)
\ee Adding (2.5.11) and (2.5.12) we have
$$
M_{ab}+(P_{cab}+P_{cba})V_c+R_{acbd}V_cV_d=0
$$
and then
$$
M_{ab}W_aW_b+(P_{cab}+P_{cba})W_aW_bV_c+R_{acbd}W_aV_cW_bV_d=0.
$$
If we write
$$
U_{ab}=\frac{1}{2}(V_aW_b-V_bW_a)=V\wedge W,
$$
then the above identity can be rearranged as \be
Q\stackrel{\Delta}{=}M_{ab}W_aW_b+2P_{abc}U_{ab}W_c
+R_{abcd}U_{ab}U_{cd}=0.%\eqno(2.5.13)
\ee This is the \textbf{Li-Yau-Hamilton quadratic}
\index{Li-Yau-Hamilton quadratic} we look for. Note that the proof
of the Li-Yau-Hamilton estimate below does not depend on the
existence of such an expanding gradient Ricci soliton. It is only
used as inspiration.

Now we are ready to state the remarkable \textbf{Li-Yau-Hamilton
estimate} \index{Li-Yau-Hamilton estimate} for the Ricci flow.

%\vskip 0.2cm\noindent {\bf Theorem 2.5.4} (Hamilton \cite{Ha93}) \
% CZ-BASELINE source-line:4492 source-kind:theorem source-id:2.5.4
\begin{theorem}[matrix Li–Yau–Hamilton estimate]
  \label{thm:cz-matrix-liyauhamilton-estimate}
  \dcref{cao-zhu:2.5.4}
  \group{Cao--Zhu Li-Yau-Hamilton Estimates}
  \source{cao-zhu}
Let $g_{ij}(x,t)$ be a complete solution with bounded curvature to
the Ricci flow on a manifold $M$ for $t$ in some time interval
$(0,T)$ and suppose the curvature operator of $g_{ij}(x,t)$ is
nonnegative. Then for any one-form $W_a$ and any two-form $U_{ab}$
we have
$$
M_{ab}W_aW_b+2P_{abc}U_{ab}W_c+R_{abcd}U_{ab}U_{cd}\geq 0
$$
on $M\times (0,T)$.
\end{theorem}

The proof of this theorem requires some rather intense
calculations. Here we only give a sketch of the proof. For more
details, we refer the reader to Hamilton's original paper
\cite{Ha93}.

\medskip
{\bf\em Sketch of the Proof.} \  Let $g_{ij}(x,t)$ be the complete
solution with bounded and nonnegative curvature operator. Recall
that in the orthonormal frame coordinate system, the curvatures
evolve by
$$
\left\{
\begin{array}{lll}
\frac{\partial}{\partial t}R_{abcd}
=\Delta R_{abcd}+2(B_{abcd}-B_{abdc}-B_{adbc} +B_{acbd}),\\[4mm]
\frac{\partial}{\partial t}R_{ab}
=\Delta R_{ab}+2R_{acbd}R_{cd},\\[4mm]
\frac{\partial}{\partial t}R=\Delta R+ 2|Ric|^2,
       \end{array}
    \right.
$$
where $B_{abcd}=R_{aebf}R_{cedf}.$

By a long but straightforward computation from these evolution
equations, one can get
$$
\(\frac{\partial}{\partial t}-\Delta\)P_{abc}=2R_{adbe}P_{dec}
+2R_{adce}P_{dbe}+2R_{bdce}P_{ade}-2R_{de}\nabla_dR_{abce}
$$
and
\begin{align*}
\(\frac{\partial}{\partial t}-\Delta\)M_{ab}
&  =2R_{acbd}M_{cd}+2R_{cd}(\nabla_cP_{dab}+\nabla_cP_{dba}) \\
&\quad +2P_{acd}P_{bcd}-4P_{acd}P_{bdc}+2R_{cd}R_{ce}R_{adbe}
-\frac{1}{2t^2}R_{ab}.
\end{align*}
Now consider
$$
Q\stackrel{\Delta}{=}M_{ab}W_aW_b+2P_{abc}U_{ab}W_c
+R_{abcd}U_{ab}U_{cd}.
$$
At a point where \be \(\frac{\partial}{\partial
t}-\Delta\)W_a=\frac{1}{t}W_a, \ \
 \(\frac{\partial}{\partial t}-\Delta\)U_{ab}=0,%\eqno(2.5.14)
\ee and \be \nabla_aW_b=0,\ \
\nabla_aU_{bc}=\frac{1}{2}(R_{ab}W_c-R_{ac}W_b)
+\frac{1}{4t}(g_{ab}W_c-g_{ac}W_b),%\eqno(2.5.15)
\ee we have
\begin{align}
\(\frac{\partial}{\partial t}-\Delta\)Q
&  = 2R_{acbd}M_{cd}W_aW_b-2P_{acd}P_{bdc}W_aW_b \\
&\quad  + 8R_{adce}P_{dbe}U_{ab}W_c+4R_{aecf}R_{bedf}U_{ab}U_{cd}\nn\\
&\quad  +
(P_{abc}W_c+R_{abcd}U_{cd})(P_{abe}W_e+R_{abef}U_{ef}).\nn
\end{align} %\eqno(2.5.16)

For simplicity we assume the manifold is compact and the curvature
operator is strictly positive. (For the general case we shall mess
the formula up a bit to sneak in the term $\epsilon e^{At}f$, as
done in Lemma 2.1.3). Suppose not; then there will be a first time
when the quantity $Q$ is zero, and a point where this happens, and
a choice of $U$ and $W$ giving the null eigenvectors. We can
extend $U$ and $W$ any way we like in space and time and still
have $Q\geq$0, up to the critical time. In particular we can make
the first derivatives in space and time to be anything we like, so
we can extend first in space to make (2.5.15) hold at that point.
And then, knowing $\Delta W_a$ and $\Delta U_{ab}$, we can extend
in time to make (2.5.14) hold at that point and that moment. Thus
we have (2.5.16) at the point.

In the RHS of (2.5.16) the quadratic term
$$
(P_{abc}W_c+R_{abcd}U_{cd})(P_{abe}W_e+R_{abef}U_{ef})
$$
is clearly nonnegative. By similar argument as in the proof of
Lemma 2.1.3, to get a contradiction we only need to show the
remaining part in the RHS of (2.5.16) is also nonnegative.

A nonnegative quadratic form can always be written as a sum of
squares of linear forms. This is equivalent to diagonalizing a
symmetric matrix and writing each nonnegative eigenvalue as a
square. Write
$$
Q=\sum_{k}(X^k_aW_a+Y^k_{ab}U_{ab})^2,\ \ \ \ \(1\leq k\leq
n+\frac{n(n-1)}{2}\).
$$
This makes
$$
M_{ab}=\sum_kX^k_aX^k_b, \ \ \ P_{abc}=\sum_kY^k_{ab}X^k_c
$$
and
$$
R_{abcd}=\sum_kY^k_{ab}Y^k_{cd}.
$$
It is then easy to compute {\allowdisplaybreaks
\begin{align*}
& 2R_{acbd}M_{cd}W_aW_b-2P_{acd}P_{bdc}W_aW_b+8R_{adce}P_{dbe}U_{ab}W_e\\
&\quad+4R_{aecf}R_{bedf}U_{ab}U_{cd}\\
&  = 2\(\sum_kY^k_{ac}Y^k_{bd}\)\(\sum_lX^l_{a}Y^l_{c}\)W_aW_b \\
&\quad-2\(\sum_kY^k_{ac}X^k_{d}\)\(\sum_lY^l_{bd}X^l_{c}\)W_aW_b\\
&\quad +8\(\sum_kY^k_{ad}Y^k_{ce}\)\(\sum_lY^l_{db}X^l_{e}\)U_{ab}W_c \\
&\quad+4\(\sum_kY^k_{ae}Y^k_{cf}\)\(\sum_lY^l_{be}Y^l_{df}\)U_{ab}U_{cd}\\
&  = \sum_{k,l}(Y^k_{ac}X^l_cW_a-Y^l_{ac}X^k_cW_a
-2Y^k_{ac}Y^l_{bc}U_{ab})^2\\
&  \geq 0.
\end{align*}
} This says that the remaining part in the RHS of (2.5.16) is also
nonnegative. Therefore we have completed the sketch of the proof.
\endproof
%%$$\eqno\#$$ \vskip 0.3cm

By taking $U_{ab}=\frac{1}{2}(V_aW_b-V_bW_a)$ and tracing over
$W_a$, we immediately get \vskip 0.2cm

%\noindent{\bf Corollary 2.5.5} (Hamilton \cite{Ha93}) \ \emph{
% CZ-BASELINE source-line:4619 source-kind:corollary source-id:2.5.5
\begin{corollary}[trace Li–Yau–Hamilton estimate]
  \label{cor:cz-trace-liyauhamilton-estimate}
  \dcref{cao-zhu:2.5.5}
  \group{Cao--Zhu Li-Yau-Hamilton Estimates}
  \source{cao-zhu}
For any one-form $V_a$ we have
$$
\frac{\partial R}{\partial t}+\frac{R}{t} +2\nabla_aR\cdot
V_a+2R_{ab}V_aV_b\geq 0.
$$
\end{corollary}
%%$$\eqno\#$$\vskip 0.3cm

In particular by taking $V\equiv 0$, we see that the function
$tR(x,t)$ is pointwise nondecreasing in time. By combining this
property with the local derivative estimate of curvature, we have
the following elliptic type estimate.

%\vskip 0.2cm \noindent {\bf Corollary 2.5.6}\ \emph {
% CZ-BASELINE source-line:4634 source-kind:corollary source-id:2.5.6
\begin{corollary}[rigidity in the trace Harnack estimate]
  \label{cor:cz-rigidity-in-the-trace-harnack-estimate}
  \dcref{cao-zhu:2.5.6}
  \group{Cao--Zhu Li-Yau-Hamilton Estimates}
  \source{cao-zhu}
Suppose we have a solution to the Ricci flow for $t>0$ which is
complete with bounded curvature, and has nonnegative curvature
operator. Suppose also that at some time $t>0$ we have the scalar
curvature $R\leq M$ for some constant M in the ball of radius $r$
around some point $p$. Then for $k=1,2,\ldots$, the $k^{th}$ order
derivatives of the curvature at $p$ at the time $t$ satisfy a
bound
$$
|\nabla^k Rm(p,t)|^2\leq
C_kM^2\(\frac{1}{r^{2k}}+\frac{1}{t^k}+M^k\)
$$
for some constant $C_k$ depending only on the dimension and $k$.
\end{corollary}

%\vskip 0.1cm \noindent {\bf Proof.}\
\begin{proof}
  \uses{thm:cz-local-curvature-derivative-estimates}
Since $tR$ is nondecreasing in time, we get a bound $R\leq 2M$ in
the given region for times between $t/2$ and $t$. The nonnegative
curvature hypothesis tells us the metric is shrinking. So we can
apply the local derivative estimate in Theorem 1.4.2 to deduce the
result.
\end{proof}
%$$\eqno\#$$ \vskip 0.3cm

By a similar argument as in Corollary 2.5.3, one readily has the
following Harnack inequality.

%\vskip 0.2cm\noindent{\bf Corollary 2.5.7} (Hamilton \cite{Ha93})\
% CZ-BASELINE source-line:4663 source-kind:corollary source-id:2.5.7
\begin{corollary}[integrated scalar-curvature Harnack estimate]
  \label{cor:cz-integrated-scalar-curvature-harnack-estimate}
  \dcref{cao-zhu:2.5.7}
  \group{Cao--Zhu Li-Yau-Hamilton Estimates}
  \source{cao-zhu}
Let $g_{ij}(x,t)$ be a complete solution of the Ricci flow on a
manifold with bounded and nonnegative curvature operator, and let
$x_1,x_2\in M,\ 0<t_1<t_2.$ Then the following inequality holds
$$
R(x_2,t_2)\geq\frac{t_1}{t_2}e^{-d_{t_1}(x_1,x_2)^2/2(t_2-t_1)}\cdot
R(x_1,t_1).$$
\end{corollary}
%$$\eqno\#$$\vskip 0.3cm

In the above discussion, we assumed that the solution to the Ricci
flow exists on $0\le t<T$, and we derived the Li-Yau-Hamilton
estimate with terms $1/t$ in it. When the solution happens to be
{\bf ancient}\index{solution!ancient}\index{ancient!solution},
i.e., defined on $-\infty<t<T$, Hamilton \cite{Ha93} found an
interesting and simple procedure for getting rid of them. Suppose
we have a solution on $\alpha<t<T$ we can replace $t$ by
$t-\alpha$ in the Li-Yau-Hamilton estimate. If we let $\alpha\to
-\infty$, then the expression $1/(t-\alpha)\to 0$ and disappears!
In particular the trace Li-Yau-Hamilton estimate in Corollary
2.5.5 becomes \be \frac{\partial R}{\partial t}+2\nabla_aR\cdot
V_a
+2R_{ab}V_aV_b\geq 0.%\eqno(2.5.17)
\ee By taking $V=0$, we see that $\frac{\partial R}{\partial
t}\geq 0$. Thus, we have the following

%\vskip 0.2cm \noindent {\bf Corollary 2.5.8} (Hamilton
%\cite{Ha93}) \ \emph{
% CZ-BASELINE source-line:4691 source-kind:corollary source-id:2.5.8
\begin{corollary}[scalar curvature monotonicity for ancient solutions]
  \label{cor:cz-scalar-curvature-monotonicity-for-ancient-solutions}
  \dcref{cao-zhu:2.5.8}
  \group{Cao--Zhu Li-Yau-Hamilton Estimates}
  \source{cao-zhu}
Let $g_{ij}(x,t)$ be a complete ancient solution of the Ricci flow
on $M\times (-\infty, T)$ with bounded and nonnegative curvature
operator, then the scalar curvature $R(x,t)$ is pointwise
nondecreasing in time $t$.
\end{corollary}
%$$\eqno\#$$\vskip 0.3cm

Corollary 2.5.8  will be very useful later on when we study
ancient $\kappa$-solutions in Chapter 6, especially combined with
Shi's derivative estimate.

We end this section by stating the Li-Yau-Hamilton estimate for
the K\"ahler-Ricci flow, due to the first author \cite{Cao92},
under the weaker curvature assumption of nonnegative holomorphic
bisectional curvature. Note that the following Li-Yau-Hamilton
estimate in the K\"ahler case is really a Li-Yau-Hamilton estimate
for the Ricci tensor of the evolving metric, so not only can we
derive an estimate on the scalar curvature, which is the trace of
the Ricci curvature, similar to Corollary 2.5.5 but also an
estimate on the determinant of the Ricci curvature as well.

%\vskip 0.2cm\noindent {\bf Theorem 2.5.9} (Cao \cite{Cao92})\
% CZ-BASELINE source-line:4714 source-kind:theorem source-id:2.5.9
\begin{theorem}[Kähler–Ricci matrix Harnack estimate]
  \label{thm:cz-kahlerricci-matrix-harnack-estimate}
  \dcref{cao-zhu:2.5.9}
  \group{Cao--Zhu Li-Yau-Hamilton Estimates}
  \source{cao-zhu}
Let $g_{\alpha\bar \beta}(x,t)$ be a complete solution to the
K\"ahler-Ricci flow on a complex  manifold $M$ with bounded
curvature and nonnegtive bisectional curvature and $0\leq t < T$.
For any point $x\in M$ and any vector $V$ in the holomorphic
tangent space $T^{1,0}_xM$, let
$$
Q_{\alpha\bar \beta}=\frac{\partial}{\partial t}R_{\alpha\bar
\beta}+ R_{\alpha\bar \gamma}R_{\gamma\bar \beta}
+\nabla_{\gamma}R_{\alpha\bar \beta} V^{\gamma}
+\nabla_{\bar\gamma}R_{\alpha\bar
\beta}V^{\bar\gamma}+R_{\alpha\bar \beta\gamma\bar \delta}V^{
\gamma}V^{\bar\delta}+\frac {1}{t}R_{\alpha\bar \beta} .
$$
Then we have
$$
Q_{\alpha\bar \beta}W^{\alpha}W^{\bar \beta}\geq 0
$$
for all $x\in M$, $V, W\in T^{1,0}_xM$, and $t>0$.
\end{theorem}
%$$\eqno\#$$

%\vskip 0.3cm \noindent{\bf Corollary 2.5.10} (Cao \cite{Cao92}) \
% CZ-BASELINE source-line:4737 source-kind:corollary source-id:2.5.10
\begin{corollary}[Kähler–Ricci scalar and determinant estimates]
  \label{cor:cz-kahlerricci-scalar-and-determinant-estimates}
  \dcref{cao-zhu:2.5.10}
  \group{Cao--Zhu Li-Yau-Hamilton Estimates}
  \source{cao-zhu}
Under the assumptions of Theorem $2.5.9,$ we have
\begin{itemize}
\item[(i)] the scalar curvature $R$ satisfies the estimate
$$
\frac{\partial R}{\partial t}-\frac{|\nabla R|^2}{R}
+\frac{R}{t}\geq 0,$$ and \item[(ii)] assuming $R_{\alpha\bar
\beta}>0$, the determinant
$\phi=\det(R_{\alpha\bar\beta})/\det(g_{\alpha\bar\beta})$ of the
Ricci curvature satisfies the estimate
$$
\frac {\partial \phi}{\partial t}-\frac{|\nabla\phi|^2}{n\phi}
+\frac {n\phi}{t}\geq 0
$$
for all $x\in M$ and $t>0$.
\end{itemize}
\end{corollary}
%$$\eqno\#$$


\section{Perelman's Estimate for Conjugate Heat Equations}

In \cite{P1} Perelman obtained a Li-Yau type estimate for
fundamental solutions of the conjugate heat equation, which is a
backward heat equation, when the metric evolves by the Ricci flow.
In this section we shall describe how to get this estimate along
the same line as in the previous section. More importantly, we
shall show how the Li-Yau path integral, when applied to
Perelman's Li-Yau type estimate, leads to an important space-time
distance function introduced by Perelman \cite{P1}. We learned
from Hamilton \cite{HaL} this idea of looking at Perelman's Li-Yau
estimate.

We saw in the previous section that the Li-Yau quantity and the
Li-Yau-Hamilton quantity vanish on expanding solutions. Note that
when we consider a backward heat equation, shrinking solitons can
be viewed as expanding backward in time. So we start by looking at
shrinking gradient Ricci solitons.

Suppose we have a shrinking gradient Ricci soliton $g_{ij}$ with
potential function $f$ on manifold $M$ and $- \infty < t < 0$ so
that, for $\tau = -t$, \be R_{ij} + \nabla_j \nabla_i f -
\frac{1}{2 \tau} g_{ij} = 0.
%\eqno{(2.6.1)}
\ee Then, by taking the trace, we have \be
R + \Delta f - \frac{n}{2 \tau} = 0. %\eqno{(2.6.2)}
\ee Also, by similar calculations as in deriving (1.1.15), we get
\be
R + |\nabla f |^2 - \frac{f}{\tau} = C %\eqno{(2.6.3)}
\ee where $C$ is a constant which we can set to be zero.

Moreover, observe \be
\frac{\partial f}{\partial t} = | \nabla f |^2 %\eqno{(2.6.4)}
\ee because $f$ evolves in time with the rate of change given by
the Lie derivative in the direction of $\nabla f$ generating the
one-parameter family of diffeomorphisms.

Combining (2.6.2) with (2.6.4), we see $f$ satisfies the backward
heat equation \be \frac{\partial f}{\partial t} = - \Delta f + |
\nabla f |^2 - R
+ \frac{n}{2 \tau}, %\eqno{(2.6.5)}
\ee or equivalently \be \frac{\partial f}{\partial \tau} = \Delta
f - | \nabla f |^2 +
R - \frac{n}{2 \tau}. %\eqno{(2.6.6)}
\ee

Recall the Li-Yau-Hamilton quadratic is a certain combination of
the second order space derivative (or first order time
derivative), first order space derivatives and zero orders.
Multiplying (2.6.2) by a factor of $2$ and subtracting (2.6.3)
yields
$$
2 \Delta f - | \nabla f |^2 + R + \frac{1}{\tau} f -
\frac{n}{\tau}=0
$$
valid for our potential function $f$ of the shrinking gradient
Ricci soliton. The quantity on the LHS of the above identity is
precisely the Li-Yau-Hamilton type quadratic found by Perelman
(cf. section 9 of \cite{P1}).

Note that a function $f$ satisfies the backward heat equation
(2.6.6) if and only if the function
$$
u=(4\pi \tau)^{-\frac{n}{2}}e^{- f}
$$
satisfies the so called \textbf{conjugate heat equation}
\index{conjugate heat equation} \be \square^\ast u \triangleq
\frac{\partial u}{\partial \tau} - \Delta
u + Ru = 0. %\eqno(2.6.7)
\ee

%\vskip 0.2cm\noindent {\bf Lemma 2.6.1} (cf. Proposition 9.1 of Perelman \cite{P1})\
% CZ-BASELINE source-line:4829 source-kind:lemma source-id:2.6.1
\begin{lemma}[gradient estimate for conjugate heat equations]
  \label{lem:cz-gradient-estimate-for-conjugate-heat-equations}
  \dcref{cao-zhu:2.6.1}
  \group{Cao--Zhu Perelman's Estimate for Conjugate Heat Equations}
  \source{cao-zhu}
Let $g_{ij}(x,t)$, $0 \leq t < T$, be a complete solution to the
Ricci flow on an $n$-dimensional manifold $M$ and let $u=(4 \pi
\tau)^{-\frac{n}{2}} e^{-f}$ be a solution to the conjugate
equation $(2.6.7)$ with $\tau=T-t$. Set
$$
H=2 \Delta f -| \nabla f |^2 + R + \frac{f-n}{\tau}
$$
and
$$
v = \tau H u = \big(\tau(R+2 \Delta f - | \nabla f|^2)+f-n \big)
u.
$$
Then we have
$$
\frac{\partial H}{\partial \tau} = \Delta H - 2 \nabla f \cdot
\nabla H - \frac{1}{\tau} H - 2 \left| R_{ij} + \nabla_i\nabla_j f
- \frac{1}{2 \tau} g_{ij} \right|^2,
$$
and
$$
\frac{\partial v}{\partial \tau} = \Delta v - Rv - 2 \tau u
\left|R_{ij}+\nabla_i\nabla_jf - \frac{1}{2 \tau} g_{ij}\right|^2.
$$
\end{lemma}

%\vskip 0.1cm \noindent {\bf Proof.}
\begin{proof}
By direct computations, we have
\begin{align*}
\frac{\partial }{\partial \tau}H &  =\frac{\partial}{\partial
\tau}\(2\triangle f -|\nabla f|^2 +R + \frac{f-n}{\tau}\)\\
&  = 2\triangle \(\frac{\partial f}{\partial \tau}\) - 2\langle
2R_{ij}, f_{ij}\rangle - 2\left\langle \nabla
f,\nabla\(\frac{\partial
f}{\partial \tau}\)\right\rangle + 2\Ric(\nabla f,\nabla f)\\
&\quad + \frac{\partial }{\partial
\tau}R+\frac{1}{\tau}\frac{\partial
}{\partial \tau}f - \frac{f-n}{\tau^2}\\
& =2\triangle \(\triangle f -|\nabla f|^2 +R -\frac{n}{2\tau}\) -
4\langle
R_{ij}, f_{ij}\rangle + 2\Ric(\nabla f,\nabla f)\\
&\quad - 2\left\langle \nabla f,\nabla\(\triangle f -|\nabla f|^2
+R
-\frac{n}{2\tau}\)\right\rangle - \triangle R -2|R_{ij}|^2\\
&\quad + \frac{1}{\tau}\(\triangle f -|\nabla f|^2 +R
-\frac{n}{2\tau}\)-\frac{f-n}{\tau^2},\\
\nabla H &  = \nabla \(2\triangle f -|\nabla f|^2 +R
+ \frac{f-n}{\tau}\)\\
&  = 2\nabla(\triangle f) - 2\langle \nabla
\nabla_if,\nabla_if\rangle + \nabla R + \frac{1}{\tau} \nabla f,\\
\triangle H &  = \triangle \(2\triangle f -|\nabla f|^2 +R
+ \frac{f-n}{\tau}\)\\
&  = 2\triangle(\triangle f) - \triangle(|\nabla f|^2)+ \triangle
R + \frac{1}{\tau} \triangle f,
\end{align*}
and
\begin{align*}
2\nabla H \cdot \nabla f &  =2\langle 2\nabla(\triangle f) -
2\langle \nabla \nabla_if,\nabla_if\rangle + \nabla R +
\frac{1}{\tau} \nabla f,\nabla f\rangle \\
&  = 2\left[\langle 2\nabla(\triangle f),\nabla f\rangle -2\langle
f_{ij},f_i f_j\rangle + \langle \nabla R,\nabla f\rangle\right] +
\frac{2}{\tau}|\nabla f|^2.
\end{align*}
Thus we get
\begin{align*}
&  \frac{\partial}{\partial \tau}H - \triangle H + 2 \nabla f
\cdot\nabla H + \frac{1}{\tau} H\\
&  =-4\langle R_{ij},f_{ij}\rangle + 2\Ric(\nabla f,\nabla f) -
2|R_{ij}|^2 -\triangle (|\nabla f|^2) + 2\langle \nabla(\triangle
f),\nabla f\rangle \\
&\quad + \frac{2}{\tau}\triangle f + \frac{2}{\tau}R - \frac{n}{2\tau^2}\\
&  = -2 \left[|R_{ij}|^2 + |f_{ij}|^2 + \frac{n}{4\tau^2} +
2\langle
R_{ij},f_{ij}\rangle -\frac{R}{\tau}-\frac{1}{\tau} \triangle f\right]\\
&  = -2\left|R_{ij} + \nabla_i\nabla_j
f-\frac{1}{2\tau}g_{ij}\right|^2,
\end{align*}
and
\begin{align*}
\(\frac{\partial}{\partial \tau} - \triangle + R\)v
& =\(\frac{\partial}{\partial \tau} - \triangle + R\)(\tau H u)\\
&  = \(\frac{\partial}{\partial \tau} - \triangle\)(\tau
H)\cdot u - 2\langle \nabla(\tau H),\nabla u\rangle\\
&  = \(\(\frac{\partial}{\partial \tau} - \triangle\)(\tau
H) - 2\langle \nabla(\tau H),\nabla f\rangle\)u\\
&  = \tau \(\frac{\partial H}{\partial \tau} - \triangle H + 2
\nabla f
\cdot \nabla H + \frac{1}{\tau} H\)u\\
&  = -2\tau u \left|R_{ij} + \nabla_i\nabla_j f
-\frac{1}{2\tau}g_{ij}\right|^2.
\end{align*}
%$$\eqno\#$$ \vskip 0.3cm
\end{proof}

Note that, since $f$ satisfies the equation (2.6.6), we can
rewrite $H$ as \be H = 2 \frac{\partial f}{\partial \tau} + |
\nabla f |^2 - R +
\frac{1}{\tau} f. %\eqno{(2.6.8)}
\ee Then, by Lemma 2.6.1, we have
$$
\frac{\partial}{\partial \tau}(\tau H)= \Delta (\tau H) -2 \nabla
f \cdot \nabla(\tau H) - 2 \tau \left|\Ric+ \nabla^2f - \frac{1}{2
\tau}g\right|^2.
$$
So by the maximum principle, we find $\operatorname{max}(\tau H)$
is nonincreasing as $\tau$ increasing.  When $u$ is chosen to be a
fundamental solution to (2.6.7), one can show that $\lim_{\tau\to
0}\tau H \leq 0$ and hence $H \leq 0$ on $M$ for all $\tau\in
(0,T]$ (see, for example, \cite{Ni06}). Since this fact is not
used in later chapters and will be only used in the rest of the
section to introduce a space-time distance via Li-Yau path
integral, we omit the details of the proof.

\medskip
Once we have Perelman's Li-Yau type estimate $H\leq 0$, we can
apply the Li-Yau path integral as in \cite{LY} to estimate the
above solution $u$ (i.e., a heat kernel estimate for the conjugate
heat equation, see also the earlier work of Cheeger-Yau
\cite{CY}). Let $p, q \in M$ be two points and $\gamma(\tau), \tau
\in [0, \bar\tau ],$ be a curve joining $p$ and $q$, with
$\gamma(0)=p$ and $\gamma(\bar{\tau})=q$. Then along the
space-time path $(\gamma(\tau), \tau)$, $\tau \in [0, \bar\tau]$,
we have
\begin{align*}
\frac{d}{d \tau} \big(2\sqrt{\tau} f (\gamma(\tau), \tau) \big) &=
2 \sqrt{\tau} \( \frac{\partial f}{\partial \tau} + \nabla
f \cdot \dot{\gamma}(\tau)\) + \frac{1}{\sqrt{\tau}}f \\
&  \leq \sqrt{\tau}\big(-|\nabla f|_{g_{ij}(\tau)}^2 + 2 \nabla
f\cdot\dot{\gamma}(\tau) \big)+\sqrt{\tau}R \\
&  = - \sqrt{\tau} |\nabla f -
\dot{\gamma}(\tau)|_{g_{ij}(\tau)}^2 +
\sqrt{\tau}(R+|\dot{\gamma}(\tau)|_{g_{ij}(\tau)}^2) \\
&  \leq \sqrt{\tau}(R + |\dot{\gamma}(\tau)|_{g_{ij}(\tau)}^2)
\end{align*}
where we have used the fact that $H\leq 0$ and the expression for
$H$ in (2.6.8).

Integrating the above inequality from $\tau =0$ to $\tau =
\bar{\tau}$, we obtain
$$
2 \sqrt{\bar{\tau}} f(q, \bar{\tau}) \leq \int^{\bar{\tau}}_0
\sqrt{\tau} (R+| \dot{\gamma}(\tau)|_{g_{ij}(\tau)}^2)d \tau,
$$

\noindent or
$$f(q, \bar{\tau}) \leq \frac{1}{2
\sqrt{\bar{\tau}}}\; \mathcal{L}(\gamma),
$$
where \be \mathcal{L}(\gamma) \triangleq \int^{\bar{\tau}}_0
\sqrt{\tau}(R+|
\dot{\gamma}(\tau)|_{g_{ij}(\tau)}^2)d\tau.  %\eqno(2.6.9)
\ee

Denote by \be l(q, \bar{\tau}) \triangleq \inf_{\gamma} \frac{1}{2
\sqrt{\bar{\tau}}}\; \mathcal{L}(\gamma),%\eqno(2.6.10)
\ee where the $inf$ is taken over all space curves $\gamma (\tau),
0 \leq \tau \leq \bar{\tau}$, joining $p$ and $q$. The space-time
distance function $l(q, \bar{\tau})$ obtained by the above Li-Yau
path integral argument is first introduced by Perelman in
\cite{P1} and is what Perelman calls reduced distance. Since
Perelman pointed out in page 19 of \cite{P1} that ``an even closer
reference in \cite{LY}, where they use `length', associated to a
linear parabolic equation, which is pretty much the same as in our
case", it is natural to call $l(q, \bar{\tau})$ the
\textbf{Li-Yau-Perelman distance}. \index{Li-Yau-Perelman
distance} See Chapter 3 for much more detailed discussions.

Finally, we conclude this section by relating the quantity $H$ (or
$v$) and the $\mathcal W$-functional of Perelman defined in
(1.5.9). Observe that $v$ happens to be the integrand of the
$\mathcal W$-functional,
$$
\mathcal{W}(g_{ij}(t), f, \tau) = \int_M v dV.
$$
Hence, when $M$ is compact,
\begin{align*}
\frac{d}{d \tau}\mathcal{W}
&  =\int_M \(\frac{\partial}{\partial \tau}v+Rv\)dV \\
& =-2\tau\int_M\left|\Ric+\nabla^2f-\frac{1}{2\tau}g\right|^2u d V\\
&  \leq 0,
\end{align*}
or equivalently,
$$
\frac{d} {d t} \mathcal{W}(g_{ij}(t),f(t),\tau(t))=\int_M
2\tau\left|R_{ij}+\nabla_i\nabla_jf-\frac
{1}{2\tau}g_{ij}\right|^2 (4\pi\tau)^{-\frac{n}{2}} e^{-f}dV,
$$
which is the same as stated in Proposition 1.5.8.

%\bigskip
